% !TeX root = ../../Book1.tex
\chapter{Sobolev 空间}
\label{b1:ch:21}
\BookChapterNumber{b1:ch:21}
% 本章摘要（不计入编号）
\begin{chapterabstract}
本章把弱导数的积分恒等式与 \(L^p\) 范数结合，建立 Sobolev 空间。先用有限组导数分量证明完备性，并在相应指数范围内证明可分性与自反性，再通过局部光滑逼近推导乘积、复合、截断及坐标变换的规则。最后选取适当的函数代表，说明一维绝对连续性怎样转成多维的沿坐标线绝对连续性，以及弱梯度怎样决定几乎处处的近似微分。空间的范数结构与代表的点态性质是两条相互联系、却不能混同的主线。
\end{chapterabstract}

% 本章目标（不计入编号）
\begin{learninggoals}
\begin{enumerate}
\item\label{b1:goal:21-definition}掌握 \(W^{k,p}\)、\(H^k\) 与局部空间的定义，区分弱导数存在、全域可积性和具体代表的正则性。
\item\label{b1:goal:21-structure}把 Sobolev 空间实现为有限乘积 \(L^p\) 的闭子空间，并据此判断完备性、可分性、自反性及分量弱收敛。
\item\label{b1:goal:21-calculus}证明乘积、链式、截断和坐标变换公式，明确局部有界性、零阶可积性以及指数端点承担的作用。
\item\label{b1:goal:21-representatives}构造一个同时适用于全部坐标方向的绝对连续代表，区分经典偏导、通常微分和近似微分，并识别后者与弱梯度的关系。
\end{enumerate}
\end{learninggoals}

本章的起点是\chapref{b1:ch:20}的弱导数及局部卷积。\secref{b1:sec:2102}使用第八至十章的 Banach 空间、Hahn--Banach 定理与 \(L^p\) 对偶。\secref{b1:sec:2103}前三类局部运算沿用第二十章的卷积；坐标变换和通常可微性还调用\cref{b1:thm:rademacher}，坐标变换另需\cref{b1:thm:weighted-area-formula}。代表理论使用\secref{b1:sec:1604}的一维绝对连续微积分与\chapref{b1:ch:14}的 Lebesgue 点；近似可微性再接上\chapref{b1:ch:13}的 Hardy--Littlewood 极大函数和\secref{b1:sec:1603}的 Lipschitz 分片。

\ProseParagraphBreak

因此，定义与空间结构只依赖前两编及卷积工具，点态代表和坐标变换还使用第三编。这些调用都发生在开集内部，既不要求边界光滑，也不预先使用下一章的全域光滑逼近。

\ProseParagraphBreak

初读时宜先贯通定义、完备性、局部运算和一维代表，再细读共同直线代表的选取与近似微分的证明。后两步中的量词尤其重要：沿几乎每条坐标线具有绝对连续性，不等于对每个代表都能逐点求导；弱梯度有界与仅属于某个有限 \(L^p\) 空间，也给出不同强度的点态结论。

\ProseParagraphBreak

Hunter 与 Nachtergaele 的《Applied Analysis》\cite[Sections 12.9--12.10]{Hunter2001Applied}提供定义和后续性质的概览，其中若干性质只陈述而不证明。需要系统伴读时，可参阅 Brézis 的《Functional Analysis, Sobolev Spaces and Partial Differential Equations》\cite{Brezis2011Functional}，或 Adams 与 Fournier 的《Sobolev Spaces》\cite{Adams2003Sobolev}。局部均值比较所用的证明路线，另在相应位置注明来源。

% !TeX root = ../../Book1.tex
\section{定义}
\label{b1:sec:2101}

上一章把导数关系写成测试积分恒等式。现在给函数及其弱导数同时加上可积性要求，就能用一个范数度量它们。这种做法适合处理取极限的问题：函数可以不再光滑，但只要各阶弱导数仍有适当的积分控制，极限就有希望留在同一个空间中。

\ProseParagraphBreak

以下固定非空开集 \(\Omega\subset\mathbb R^d\)，使用 Lebesgue 测度，标量域为 \(\mathbb K=\mathbb R\) 或 \(\mathbb C\)。不预设 \(\Omega\) 有界、连通或边界光滑。弱导数始终按\cref{b1:def:weak-derivative}理解；函数及其导数均按几乎处处相等取等价类。
% 实读来源：Hunter2001Applied作者公开第12章印刷362--364页，定义12.66及其后的范数、内积。
% 此处完整规定本书多指标范数；对局部空间、边界可积性及高阶计数另作说明。

\subsection{\texorpdfstring{\(W^{1,p}\)}{W1,p}}
\label{b1:subsec:210101}

\begin{definition}\label{b1:def:sobolev-first-order}
给定 \(1\le p\le\infty\)，若 \(u\in L^p(\Omega)\)，且它的每个一阶弱偏导数 \(\partial_i u\) 均存在并属于 \(L^p(\Omega)\)，则称 \(u\) 属于一阶\term[sobolev-kongjian]{Sobolev 空间}[Sobolev space]，记为 \(u\in W^{1,p}(\Omega)\)。规定
\[
 \|u\|_{W^{1,p}(\Omega)}
 \coloneqq\left(\|u\|_{L^p(\Omega)}^p+
         \sum_{i=1}^d\|\partial_i u\|_{L^p(\Omega)}^p\right)^{1/p}
 \quad(1\le p<\infty),
\]
\[
 \|u\|_{W^{1,\infty}(\Omega)}
 \coloneqq\max\bigl\{\|u\|_{L^\infty(\Omega)},
                \|\partial_1u\|_{L^\infty(\Omega)},\ldots,
                \|\partial_du\|_{L^\infty(\Omega)}\bigr\}.
\]
\end{definition}
\symidx{\(W^{1,p}(\Omega)\)：函数及一阶弱导数属于\(L^p\)的空间}

这些弱导数首先确实可以按上一章定义。因为对每个紧集 \(K\Subset\Omega\)，Hölder 不等式给出
\[
 \int_K|u|\,\mathrm dx
 \le m_d(K)^{1-1/p}\|u\|_{L^p(\Omega)}
 \quad(1\le p<\infty),
\]
而 \(p=\infty\) 时右端换为 \(m_d(K)\|u\|_\infty\)。因此 \(L^p(\Omega)\subset L^1_{\mathrm{loc}}(\Omega)\)，即使 \(\Omega\) 的总体积无限也成立。\secref{b1:sec:2004}的唯一性又保证范数不依赖弱导数代表的选择。

\ProseParagraphBreak

定义中既控制导数，也控制函数本身。在有界开集上，非零常函数的梯度为零，函数本身却不是零等价类；因此只用 \(\|\nabla u\|_p\) 不能定义这里的范数。涉及边界条件或把常数取商的空间，需要以后另作定义。

\ProseParagraphBreak

在 \((-1,1)\) 上，\(u(x)=|x|\) 属于每个 \(W^{1,p}\)，包括 \(p=\infty\)，尽管它在原点不可微。相反，一个函数即使在开集内处处光滑，也不保证满足全域积分条件。例如 \(u(x)=\sqrt{x}\) 在 \((0,1)\) 内光滑，且函数本身有界，但
\[
 \int_0^1|u'(x)|^p\,\mathrm dx
 =2^{-p}\int_0^1x^{-p/2}\,\mathrm dx
\]
只在 \(p<2\) 时有限。因此它属于 \(W^{1,p}(0,1)\) 当且仅当 \(1\le p<2\)。这里的障碍发生在趋近边界时，不在任何内部点。

\subsection{\texorpdfstring{\(W^{k,p}\)}{Wk,p}}
\label{b1:subsec:210102}

令 \(k\in\mathbb N_0\)，记
\[
 A_k\coloneqq\{\alpha\in\mathbb N_0^d:|\alpha|\le k\},
 \qquad N_k\coloneqq\#A_k=\binom{d+k}{k}.
\]
约定 \(\partial^0u=u\)。对高阶空间，需要每一个不超过 \(k\) 阶的弱导数，而不只是某几个最高阶组合。

\begin{definition}\label{b1:def:sobolev-higher-order}
若 \(u\in L^p(\Omega)\)，且对每个 \(\alpha\in A_k\)，弱导数 \(\partial^\alpha u\) 都属于 \(L^p(\Omega)\)，则称 \(u\in W^{k,p}(\Omega)\)。其范数规定为
\begin{equation}\label{b1:eq:sobolev-norm}
 \|u\|_{W^{k,p}(\Omega)}
 \coloneqq\left(\sum_{\alpha\in A_k}
              \|\partial^\alpha u\|_{L^p(\Omega)}^p\right)^{1/p}
 \quad(p<\infty),
\end{equation}
而 \(p=\infty\) 时取这些导数的 \(L^\infty\) 范数的最大值。特别地，\(W^{0,p}(\Omega)=L^p(\Omega)\)。
\end{definition}
\symidx{\(W^{k,p}(\Omega)\)：函数及所有不超过\(k\)阶的弱导数均属于\(L^p\)的空间}

对整数 \(m\geq1\)，正式约定
\[
 W^{k,p}(\Omega;\mathbb R^m)
 \coloneqq\bigl\{u=(u^1,\ldots,u^m):u^a\in W^{k,p}(\Omega),\ 1\leq a\leq m\bigr\}.
\]
这里
\[
 \partial^\alpha u
 \coloneqq\bigl(\partial^\alpha u^1,\ldots,\partial^\alpha u^m\bigr),
\]
并固定用 \(\mathbb R^m\) 的欧氏范数计算每个 \(\|\partial^\alpha u\|_{L^p(\Omega;\mathbb R^m)}\)。向量值 Sobolev 范数仍按\eqref{b1:eq:sobolev-norm}对 \(\alpha\in A_k\) 作 \(p\) 次和；当 \(p=\infty\) 时，取这些向量值 \(L^\infty\) 范数的最大值。这一约定固定了后文定量估计中的有限维范数，而不只依赖有限维范数等价性。
\symidx{\(W^{k,p}(\Omega;\mathbb R^m)\)：各分量属于\(W^{k,p}(\Omega)\)的有限维向量值 Sobolev 空间}

\ProseParagraphBreak

若 \(|\alpha|+|\beta|\le k\)，则 \(\partial^\alpha(\partial^\beta u)=\partial^{\alpha+\beta}u\)。这不是另加一条经典可微假设：等式先由分布导数可交换得到，再由局部可积代表的唯一性成为弱导数的等式。因此也可以递归地检查 \(u\) 及其低阶导数是否具有下一阶弱导数。

\ProseParagraphBreak

同一空间常配有不同但等价的范数。例如令
\[
 S(u)=\sum_{\alpha\in A_k}\|\partial^\alpha u\|_p.
\]
对 \(p<\infty\)，有限维 Hölder 不等式给出
\[
 \|u\|_{W^{k,p}}\le S(u)
 \le N_k^{1-1/p}\|u\|_{W^{k,p}},
\]
对 \(p=\infty\) 则有 \(\|u\|_{W^{k,\infty}}\le S(u)\le N_k\|u\|_{W^{k,\infty}}\)。所以空间及收敛概念相同，但后续若要追踪精确常数，须回到本书的\eqref{b1:eq:sobolev-norm}，不能把等价范数当作数值相等。

\ProseParagraphBreak

当 \(k\ge1\) 时，记最高阶部分为
\[
 |u|_{W^{k,p}(\Omega)}
 \coloneqq\left(\sum_{|\alpha|=k}\|\partial^\alpha u\|_p^p\right)^{1/p}
 \quad(p<\infty),
\]
无穷端点仍取最大值。这一般只是半范数。比如在有界域上，次数低于 \(k\) 的多项式具有零的最高阶半范数，却未必为零函数。

\ProseParagraphBreak

提高 \(k\) 会缩小空间：\(W^{k+1,p}(\Omega)\subset W^{k,p}(\Omega)\)，且本书范数下包含映射的范数不超过 \(1\)。提高 \(p\) 则还要看域的测度。若 \(m_d(\Omega)<\infty\) 且 \(1\le p\le q\le\infty\)，逐项使用 Hölder 可得
\[
 \|u\|_{W^{k,p}(\Omega)}
 \le N_k^{1/p-1/q}\cdot m_d(\Omega)^{1/p-1/q}
       \|u\|_{W^{k,q}(\Omega)}.
\]
在无限测度域上没有这条无条件包含关系；常函数 \(1\) 在 \(\mathbb R^d\) 上属于 \(W^{k,\infty}\)，却不属于任何有限 \(p\) 的 \(W^{k,p}\)。

\subsection{\texorpdfstring{\(H^k=W^{k,2}\)}{Hk=Wk,2}}
\label{b1:subsec:210103}

\begin{definition}\label{b1:def:sobolev-hilbert}
将 \(W^{k,2}(\Omega)\) 记作 \(H^k(\Omega)\)，并定义
\begin{equation}\label{b1:eq:sobolev-inner-product}
 \langle u,v\rangle_{H^k(\Omega)}
 \coloneqq\sum_{\alpha\in A_k}\int_\Omega
       \partial^\alpha u\,\overline{\partial^\alpha v}\,\mathrm dx.
\end{equation}
实值情形省略复共轭。
\end{definition}
\symidx{\(H^k(\Omega)\)：\(W^{k,2}(\Omega)\)及其导数内积}

内积对第一变量线性，与\chapref{b1:ch:08}约定一致。零阶项保证正定性，并且 \(\langle u,u\rangle_{H^k}=\|u\|_{W^{k,2}}^2\)。因此 \(H^k\) 不仅表示另一种记号，还突出了二次能量与正交方法可以使用的结构；要称它为 Hilbert 空间，仍须在下一节证明完备性。

\ProseParagraphBreak

这里每个多重指标只计一次。例如二维 \(H^2\) 范数含有一个 \(\|\partial_1\partial_2u\|_2^2\) 项；若把 Hessian 矩阵每个位置都逐项平方求和，混合导数会计两次，得到的是另一个等价范数。两种约定都常用，本文需要数值恒等式时始终以\eqref{b1:eq:sobolev-inner-product}为准。

\ProseParagraphBreak

Hunter 与 Nachtergaele 的《Applied Analysis》\cite[Section 12.9, Definition 12.66 and the following discussion]{Hunter2001Applied}把这些空间放在弱导数之后引入。接下来我们将把导数作为 \(L^p\) 中的有限组分量，直接从\chapref{b1:ch:10}的完备性推导其基本结构。

\subsection{局部 Sobolev 空间}
\label{b1:subsec:210104}

全域可积性会同时检测函数在边界附近及无穷远处的行为。研究开集内部的运算时，往往只需要把这些位置暂时排除。

\begin{definition}\label{b1:def:sobolev-local}
若 \(u\in L^1_{\mathrm{loc}}(\Omega)\)，且每个 \(|\alpha|\le k\) 阶弱导数都存在并属于 \(L^p_{\mathrm{loc}}(\Omega)\)，则称 \(u\) 属于\term[jubu-sobolev-kongjian]{局部 Sobolev 空间}[local Sobolev space]，记作 \(W^{k,p}_{\mathrm{loc}}(\Omega)\)。这里 \(L^p_{\mathrm{loc}}\) 表示在每个紧子集上 \(p\) 次幂可积，或在 \(p=\infty\) 时本性有界。
\end{definition}
\symidx{\(W^{k,p}_{\mathrm{loc}}(\Omega)\)：局部Sobolev空间}

向量值局部空间也逐分量定义：
\[
 W^{k,p}_{\mathrm{loc}}(\Omega;\mathbb R^m)
 \coloneqq\bigl\{u=(u^1,\ldots,u^m):u^a\in W^{k,p}_{\mathrm{loc}}(\Omega),\ 1\leq a\leq m\bigr\}.
\]
因此后文的向量值链式法则不再把“逐分量理解”作为未写出的定义。

\ProseParagraphBreak

等价地，对每个开集 \(V\) 满足 \(\overline V\Subset\Omega\)，都有 \(u|_V\in W^{k,p}(V)\)。一个方向由限制立即得到。反过来，用可数个这类 \(V\) 覆盖 \(\Omega\)；重叠部分的弱导数由唯一性相等，故能拼接为同一个局部 \(L^p\) 函数，再用\cref{b1:prop:weak-derivative-locality}得到整个 \(\Omega\) 上的弱导数关系。

\ProseParagraphBreak

例如 \(x\mapsto1/x\) 属于 \((0,1)\) 上的每个 \(W^{k,p}_{\mathrm{loc}}\)，但并不属于该区间上的 \(L^1\)。另一个方向，\(W^{k,p}(\Omega)\) 总包含于 \(W^{k,p}_{\mathrm{loc}}(\Omega)\)，而局部成员只有在所有导数的全域范数也有限时才回到全域空间。

\ProseParagraphBreak

局部收敛同样按内部区域理解：\(u_j\to u\) 于 \(W^{k,p}_{\mathrm{loc}}(\Omega)\)，是指对上述每个 \(V\)，都有 \(\|u_j-u\|_{W^{k,p}(V)}\to0\)。只要选一列内部开集穷竭，使每个紧子集最终落在其中，检查这可数列就足够。局部化并不规定边界值，更不意味着把函数在域外补零后仍属于同阶 Sobolev 空间；上一章的端点点质量已经说明零延拓可能产生额外的分布导数。

% !TeX root = ../../Book1.tex
\section{基本结构}
\label{b1:sec:2102}

弱导数的定义把函数与一组积分恒等式联系起来，Sobolev 范数则同时控制函数和这些导数。两者配合，才能回答极限是否仍留在空间内。本节把所有导数放入一个有限乘积空间：范数、完备性和弱收敛都可以在这些分量上研究，而分量之间必须满足的相容关系由测试函数保存。

\ProseParagraphBreak

Hunter 与 Nachtergaele 的《Applied Analysis》\cite[Section 12.9, p. 363]{Hunter2001Applied}在 Sobolev 空间定义之后列出 Banach 与 Hilbert 结构。本节从\chapref{b1:ch:10}的 \(L^p\) 理论出发证明这些结论，并展开连续泛函如何由导数分量表示。以下固定非空开集 \(\Omega\subseteq\mathbb R^d\)、\(k\in\mathbb N_0\) 及标量域 \(\mathbb K=\mathbb R\) 或 \(\mathbb C\)，不要求 \(\Omega\) 有界或具有规则边界。

% 实读来源：Hunter2001Applied，作者官方 ch12.pdf，印刷 pp.362--367；
% 定义12.66及Banach/Hilbert结构的陈述在p.363，未附证明。
% pp.365--367为作者明确不附证明的性质综述，不作为本节技术证明来源。
% 下文有限乘积、闭性、可分性及对偶证明以本书第8--10章已证结果展开，
% 不使用第22章的全局光滑稠密性或边界延拓。

\subsection{导数分量与完备性}
\label{b1:subsec:210201}

沿用前节的有限指标集 \(A_k=\{\,\alpha\in\mathbb N_0^d\mid|\alpha|\leq k\,\}\)，并置
\[
 Y_p=\prod_{\alpha\in A_k}L^p(\Omega;\mathbb K).
\]
对 \(v=(v_\alpha)_{\alpha\in A_k}\)，赋予范数
\[
 \|v\|_{Y_p}=
 \begin{cases}
  \left(\displaystyle\sum_{\alpha\in A_k}\|v_\alpha\|_p^p\right)^{1/p},
       &1\leq p<\infty,\\[2mm]
  \displaystyle\max_{\alpha\in A_k}\|v_\alpha\|_\infty,
       &p=\infty.
 \end{cases}
\]
这里各分量相互独立；任取这样的函数族，通常不能把它们全部看成同一个函数的导数。

\begin{proposition}\label{b1:prop:sobolev-derivative-embedding}
映射
\[
 \mathcal J:W^{k,p}(\Omega)\longrightarrow Y_p,
 \quad u\longmapsto(\partial^\alpha u)_{\alpha\in A_k}
\]
是线性等距嵌入。特别地，前一节规定的 \(\|\cdot\|_{W^{k,p}}\) 确为范数，每个导数映射
\(u\mapsto\partial^\alpha u\) 都是范数不超过 \(1\) 的有界线性映射。
\end{proposition}
\begin{proof}
弱导数恒等式对函数线性，唯一性保证这些运算在几乎处处等价类上有定义。有限 \(p\) 的三角不等式来自各分量中的 Minkowski 不等式及有限序列的 Minkowski 不等式；后者是对有限计数测度应用\cref{b1:thm:lp-minkowski}。无穷端点直接取各分量范数的最大值。范数等式
\(\|\mathcal Ju\|_{Y_p}=\|u\|_{W^{k,p}}\) 来自定义。由于指标集含零指标，\(\mathcal Ju=0\) 特别蕴含 \(u=0\) 几乎处处，所以映射单射，范数也正定。最后，每个分量的范数不超过整个乘积范数，得到导数映射的估计。
\end{proof}

\begin{theorem}[Sobolev 空间的完备性]\label{b1:thm:sobolev-complete}
对每个 \(1\leq p\leq\infty\)，像空间 \(\mathcal J\bigl(W^{k,p}(\Omega)\bigr)\) 是 \(Y_p\) 的闭线性子空间，因而 \(W^{k,p}(\Omega)\) 是 Banach 空间。
\end{theorem}
\begin{proof}
先看 \(Y_p\)。其中的 Cauchy 列在每个分量上都是 \(L^p\) 中的 Cauchy 列。由\cref{b1:thm:lp-complete}，各分量都有极限；分量只有有限个，所以这些极限组成的向量也是原序列在 \(Y_p\) 中的极限。因此 \(Y_p\) 完备。

设 \(u_j\in W^{k,p}(\Omega)\)，并且 \(\mathcal Ju_j\to(v_\alpha)_\alpha\) 于 \(Y_p\)。取零阶分量 \(u=v_0\)。我们需要证明 \(v_\alpha=\partial^\alpha u\)，不能只把它们命名为导数。给定 \(\varphi\in C_c^\infty(\Omega)\)，记 \(p'\) 为共轭指数，端点取 \(1'=\infty\)、\(\infty'=1\)。函数 \(\varphi\) 及其各阶导数均有紧支集，故属于 \(L^{p'}(\Omega)\)。Hölder 不等式给出
\[
 \left|\int_\Omega(\partial^\alpha u_j-v_\alpha)\varphi\,\mathrm dx\right|
 \leq\|\partial^\alpha u_j-v_\alpha\|_p\cdot\|\varphi\|_{p'}
 \longrightarrow0,
\]
也给出 \(\int_\Omega u_j\partial^\alpha\varphi\to\int_\Omega u\partial^\alpha\varphi\)。于是对每个 \(\alpha\in A_k\)，
\[
 \int_\Omega v_\alpha\varphi\,\mathrm dx
 =\lim_{j\to\infty}\int_\Omega\partial^\alpha u_j\,\varphi\,\mathrm dx
 =(-1)^{|\alpha|}\int_\Omega u\partial^\alpha\varphi\,\mathrm dx.
\]
所有分量均局部可积，所以上式正是弱导数的定义。因而 \(u\in W^{k,p}(\Omega)\) 且 \(\mathcal Ju=(v_\alpha)_\alpha\)，像空间闭。最后应用\cref{b1:prop:closed-subspace-complete}，再通过等距映射 \(\mathcal J\) 把完备性传回原空间。
\end{proof}

证明中没有在 \(\partial\Omega\) 上做任何运算。测试函数的支集留在域内，使极限所需的积分始终受控；这也是任意开集都适用的原因。另一方面，不能在乘积空间中随意指定导数。例如在 \((0,1)\) 上，\((0,1)\in L^p\times L^p\) 不可能等于 \((u,u')\)：零阶分量要求 \(u=0\)，其弱导数便也只能为零。

\subsection{Hilbert 结构}
\label{b1:subsec:210202}

\begin{corollary}\label{b1:cor:sobolev-hilbert}
空间 \(H^k(\Omega)=W^{k,2}(\Omega)\) 关于内积
\[
 \langle u,v\rangle_{H^k}
 =\sum_{\alpha\in A_k}
   \int_\Omega\partial^\alpha u\,\overline{\partial^\alpha v}\,\mathrm dx
\]
是 Hilbert 空间，其内积范数恰为前一节的 \(W^{k,2}\) 范数。
\end{corollary}
\begin{proof}
每一项由 \(L^2\) 中的 Cauchy--Schwarz 不等式保证有限。积分的线性性与复共轭规则给出第一变量线性和共轭对称性；且
\[
 \langle u,u\rangle_{H^k}
 =\sum_{\alpha\in A_k}\|\partial^\alpha u\|_2^2
 =\|u\|_{W^{k,2}}^2.
\]
零阶项保证正定性。前一定理已证明该范数完备，因此它是 Hilbert 内积。
\end{proof}

这使\chapref{b1:ch:08}的正交投影与表示定理可以用于 Sobolev 函数。具体地，由\cref{b1:thm:hilbert-riesz}，每个 \(\Lambda\in H^k(\Omega)^*\) 都有唯一 \(w\in H^k(\Omega)\)，使
\[
 \Lambda(u)=\sum_{\alpha\in A_k}
       \int_\Omega\partial^\alpha u\,\overline{\partial^\alpha w}\,\mathrm dx,
 \quad \|\Lambda\|=\|w\|_{H^k}.
\]
复情形中，\(w\mapsto\Lambda\) 是共轭线性的，不能把上式中的共轭静默去掉。

\ProseParagraphBreak

例如，给定这个泛函，考虑
\[
 \mathcal E(u)=\dfrac12\|u\|_{H^k}^2-\operatorname{Re}\Lambda(u).
\]
把 \(\Lambda(u)=\langle u,w\rangle_{H^k}\) 代入并展开内积，得到
\[
 \mathcal E(u)=\dfrac12\|u-w\|_{H^k}^2-\dfrac12\|w\|_{H^k}^2.
\]
因此最小值在 \(u=w\) 处唯一达到。这里同时控制了零阶及高阶导数；若删去零阶项，常数函数可能有零能量，便不能不加条件地沿用同一范数和唯一性论证。

\subsection{可分性与自反性}
\label{b1:subsec:210203}

\begin{proposition}\label{b1:prop:sobolev-separable-reflexive}
若 \(1\leq p<\infty\)，则 \(W^{k,p}(\Omega)\) 可分；若 \(1<p<\infty\)，则它还自反。
\end{proposition}
\begin{proof}
先给出 \(L^p(\Omega)\) 的可数稠密集。由\cref{b1:prop:lp-continuous-dense}，只须在 \(\mathbb R^d\) 上逼近连续紧支集函数。把空间划成边长 \(2^{-j}\) 的半开立方体，在有限多个这样的立方体上取有理数值，其余处取零；复情形取实部、虚部均有理的数。所有这些阶梯函数组成可数集。

给定 \(f\in C_c(\mathbb R^d)\)，固定一个内部包含其支集的大立方体。对充分细的网格，仅在与支集相交的网格立方体上取样，再把样本值逼近为上述有理数值。函数 \(f\) 一致连续，故所得阶梯函数与 \(f\) 的一致误差趋于零；误差支集又全落在一个固定有界集内，因此 \(L^p\) 误差也趋于零。任意 \(L^p(\Omega)\) 函数在域外补零后属于 \(L^p(\mathbb R^d)\)，将上述逼近限制回 \(\Omega\)，便得到所需可数稠密集。这里的补零只用于 \(L^p\) 函数，没有断言它保持弱导数。

有限乘积 \(Y_p\) 因而可分。可分度量空间的子空间仍可分：用一个可数稠密集中的点作中心、正有理数作半径，得到可数拓扑基；对其中每个与子空间相交的基元素，选取一个交点，这些交点便在子空间中稠密。将此事实用于 \(\mathcal J\bigl(W^{k,p}(\Omega)\bigr)\)，可分性得证。

为证明自反性，把有限乘积看成一个 \(L^p\) 空间。令
\[
 Z=\Omega\times A_k,
 \quad \nu(A)=\sum_{\alpha\in A_k}
       m_d\bigl(\{\,x\in\Omega\mid(x,\alpha)\in A\,\}\bigr),
\]
其中可测集是各截面都为 Lebesgue 可测的集合。有限求和定义了一个 \(\sigma\)-有限测度，映射 \((v_\alpha)_\alpha\mapsto[(x,\alpha)\mapsto v_\alpha(x)]\) 将 \(Y_p\) 线性等距地识别为 \(L^p(\nu)\)。若 \(1<p<\infty\)，由\cref{b1:thm:lp-reflexive}，\(Y_p\) 自反；其闭线性子空间由\cref{b1:cor:closed-subspace-reflexive}也自反。通过 \(\mathcal J\) 即得结论。
\end{proof}

有限 \(p\) 的可分性并不来自“所有 Sobolev 函数都能由光滑函数逼近”这一尚未证明的命题。我们只使用了 \(L^p\) 的可分性及度量子空间的性质；由此选出的稠密族未必光滑。

\ProseParagraphBreak

无穷端点不能作同样结论，甚至 \(W^{1,\infty}(0,1)\) 也不可分。函数本身可以连续且一致有界，其导数中的跳跃位置仍能在无穷范数下彼此分离；本章习题将给出这一边界的具体构造。

\subsection{弱收敛}
\label{b1:subsec:210204}

范数收敛等价于有限多个导数分量的 \(L^p\) 范数收敛。弱收敛也有相应的分量刻画，但不能仅由范数定义直接宣布这一点：还需说明原空间上的每个连续泛函都能由分量上的泛函组合而成。

\begin{proposition}\label{b1:prop:sobolev-functional-representation}
设 \(1\leq p<\infty\)，\(q=p'\)。每个 \(\Lambda\in W^{k,p}(\Omega)^*\) 都可写为
\begin{equation}\label{b1:eq:sobolev-functional-representation}
 \Lambda(u)=\sum_{\alpha\in A_k}
   \int_\Omega g_\alpha\,\partial^\alpha u\,\mathrm dx,
 \quad g_\alpha\in L^q(\Omega).
\end{equation}
反过来，每个这样的有限函数族都定义连续线性泛函，且
\[
 \|\Lambda\|\leq\|(g_\alpha)_\alpha\|_{Y_q}.
\]
表示族通常不唯一，但可选得此处等号成立；换言之，\(\|\Lambda\|\) 是所有表示族的 \(Y_q\) 范数的最小值。
\end{proposition}
\begin{proof}
在 \(\mathcal J\bigl(W^{k,p}(\Omega)\bigr)\) 上定义 \(F(\mathcal Ju)=\Lambda(u)\)。等距性保证 \(F\) 良定义且范数等于 \(\|\Lambda\|\)。由\cref{b1:thm:hahn-banach}，它有同范数延拓 \(\widetilde F\in Y_p^*\)。使用上一证明中的 \(\sigma\)-有限测度空间 \((Z,\nu)\)，由\cref{b1:thm:lp-dual}及 \(p=1\) 时的\cref{b1:thm:l1-dual-sigma-finite}，存在 \(g\in L^q(\nu)\)，满足
\[
 \widetilde F(v)=\int_Zgv\,\mathrm d\nu,
 \quad\|g\|_{L^q(\nu)}=\|\widetilde F\|=\|\Lambda\|.
\]
写成截面 \(g_\alpha(x)=g(x,\alpha)\)，再限制到 \(v=\mathcal Ju\)，即得所需表示和等范数。反方向在 \((Z,\nu)\) 上应用 Hölder 不等式，得到
\[
 |\Lambda(u)|
 \leq\|(g_\alpha)_\alpha\|_{Y_q}\cdot\|\mathcal Ju\|_{Y_p}
 =\|(g_\alpha)_\alpha\|_{Y_q}\cdot\|u\|_{W^{k,p}}.
\]
于是每个表示族的范数均不小于 \(\|\Lambda\|\)，先前构造的族达到此下界。
\end{proof}

这个表示沿用 \(L^p\) 对偶的双线性积分约定，\(g_\alpha\) 上不取共轭。它与 Hilbert 表示的区别还在于各分量未被要求来自同一函数。非唯一性可以直接看见：若 \(k\geq1\)、\(\psi\in C_c^\infty(\Omega)\)，取 \(g_0=\partial_1\psi\)、\(g_{e_1}=\psi\)，其余分量取零，则弱导数恒等式给出
\[
 \int_\Omega u\partial_1\psi\,\mathrm dx
 +\int_\Omega\psi\partial_1u\,\mathrm dx=0.
\]
所以可以把这一族加到任何表示族上而不改变泛函。它并不抵触 Hilbert 表示向量 \(w\) 的唯一性。

\begin{proposition}\label{b1:prop:sobolev-weak-components}
设 \(1\leq p\leq\infty\)，\(u_j,u\in W^{k,p}(\Omega)\)。则
\[
 u_j\rightharpoonup u\text{ 于 }W^{k,p}(\Omega)
 \quad\Longleftrightarrow\quad
 \partial^\alpha u_j\rightharpoonup\partial^\alpha u
 \text{ 于 }L^p(\Omega)\quad(\alpha\in A_k).
\]
当 \(p<\infty\) 时，右边又等价于对每个 \(\alpha\in A_k\) 和 \(g\in L^{p'}(\Omega)\)，
\[
 \int_\Omega g\partial^\alpha u_j\,\mathrm dx
 \longrightarrow\int_\Omega g\partial^\alpha u\,\mathrm dx.
\]
\end{proposition}
\begin{proof}
若在 Sobolev 空间中弱收敛，将 \(L^p\) 上任意连续线性泛函与有界导数映射 \(u\mapsto\partial^\alpha u\) 复合，便得到各分量的弱收敛。

反过来，设每个分量弱收敛。对任意 \(\Lambda\in W^{k,p}(\Omega)^*\)，再次用 Hahn--Banach 定理将其通过 \(\mathcal J\) 延拓为 \(\widetilde F\in Y_p^*\)。记 \(\iota_\alpha:L^p(\Omega)\to Y_p\) 为只在第 \(\alpha\) 个分量放入输入函数的映射，令 \(F_\alpha=\widetilde F\circ\iota_\alpha\)。这些都是 \(L^p\) 上的连续线性泛函，且有限分量分解给出
\[
 \Lambda(u_j)=\sum_{\alpha\in A_k}F_\alpha(\partial^\alpha u_j)
 \longrightarrow\sum_{\alpha\in A_k}F_\alpha(\partial^\alpha u)
 =\Lambda(u).
\]
此论证也适用于 \(p=\infty\)。有限 \(p\) 时再用 \(L^p\) 对偶表示，便得到积分形式的刻画。
\end{proof}

在无穷端点，只用 \(L^1\) 函数检验各分量，表达的是各分量的\WeakStar{}收敛，不能据此改称 \(W^{k,\infty}\) 的弱收敛。上一个命题的端点版本使用整个 \((L^\infty)^*\)。

\begin{corollary}\label{b1:cor:sobolev-weak-subsequence}
若 \(1<p<\infty\) 且 \((u_j)\) 在 \(W^{k,p}(\Omega)\) 中有界，则存在子列 \((u_{j_\ell})\) 及 \(u\in W^{k,p}(\Omega)\)，使 \(u_{j_\ell}\rightharpoonup u\)。并且
\[
 \|u\|_{W^{k,p}}\leq\liminf_{\ell\to\infty}\|u_{j_\ell}\|_{W^{k,p}}.
\]
\end{corollary}
\begin{proof}
由自反性及\cref{b1:thm:reflexive-weak-subsequence}取得弱收敛子列；最后的不等式来自\cref{b1:prop:norm-weak-lsc}。
\end{proof}

也可以先对各阶导数应用 \(L^p\) 的弱子列结论。分量有限，逐次抽取后有共同子列；把测试函数代入各分量的弱极限，如完备性证明那样即可辨认零阶极限的弱导数。前一命题保证最后得到的确实是 Sobolev 空间中的弱收敛。这种分量做法常用于逐阶估计近似解。

\ProseParagraphBreak

当 \(p=1\) 时，有界性本身不再保证上述抽列结论，导数质量可能集中在越来越小的区域内。阻止这种集中的条件属于\secref{b1:sec:1005}的一致可积性理论，不由 Sobolev 范数的有界性自动产生。章末将用一维有界序列说明这个障碍。

% !TeX root = ../../Book1.tex
\section{运算规则}
\label{b1:sec:2103}

经典微积分的运算公式能否保留，要同时回答两个问题：等式在分布意义下是否成立，右端是否具有所要求的可积性。下面先在开集内部把函数光滑化，再用局部积分收敛传递公式。整个过程只使用第20章的卷积，不依赖下一章的全域光滑稠密性或延拓定理。

\begin{lemma}\label{b1:lem:sobolev-local-smoothing}
设 \(u\in W^{k,p}_{\mathrm{loc}}(\Omega)\)，\(U,V\) 为开集，\(\overline U\Subset V\)、\(\overline V\Subset\Omega\)。取第20章的非负单位积分磨光函数，令 \(u_\varepsilon=u*\rho_\varepsilon\)。当 \(\varepsilon\) 充分小时，\(u_\varepsilon\) 在 \(U\) 的邻域光滑，且
\[
 \partial^\alpha u_\varepsilon
   =(\partial^\alpha u)*\rho_\varepsilon
 \quad\text{在 }U\text{ 上},\qquad |\alpha|\le k.
\]
若 \(p<\infty\)，则 \(u_\varepsilon\to u\) 于 \(W^{k,p}(U)\)。若 \(p=\infty\)，则对每个有限 \(q\ge1\) 有 \(W^{k,q}(U)\) 收敛，并有
\[
 \|u_\varepsilon\|_{W^{k,\infty}(U)}
 \le\|u\|_{W^{k,\infty}(V)}.
\]
\end{lemma}
\begin{proof}
选择 \(\varepsilon<\operatorname{dist}(\overline U,V^c)\)。卷积在 \(U\) 上只读取 \(V\) 中的值，导数交换由\cref{b1:prop:distribution-convolution-smooth}成立。对有限个导数分别应用\cref{b1:prop:local-smoothing-approximation}，得到有限 \(p\) 的范数收敛。若 \(p=\infty\)，这些导数在 \(V\) 上也属于每个有限 \(L^q\)，故局部收敛仍成立；非负核与单位积分则给出每个导数的本性界，取最大值即得所列估计。
\end{proof}

后面的证明还会使用这样的取极限方式：若 \(u_j\to u\)、\(g_{i,j}\to g_i\) 于 \(L^1_{\mathrm{loc}}\)，且 \(\partial_i u_j=g_{i,j}\)，则 \(\partial_i u=g_i\)。只需对固定测试函数在分部积分恒等式两端取极限，正如第\ref{b1:sec:2004}节已说明的那样。

无穷端点只提供有限 \(q\) 的局部逼近和统一界，不能改写为 \(W^{k,\infty}\) 范数逼近。例如 \(|x|\) 的弱导数为符号函数，光滑函数不可能在原点附近依本性上确界范数逼近这个跳跃。

\subsection{乘积法则}
\label{b1:subsec:210301}

先考虑局部有界的因子，这使乘积右端仍处于原来的 \(L^p\) 尺度。

\begin{proposition}[Leibniz]\label{b1:prop:sobolev-product}
若 \(u,v\in W^{1,p}_{\mathrm{loc}}(\Omega)\cap L^\infty_{\mathrm{loc}}(\Omega)\)，则 \(uv\in W^{1,p}_{\mathrm{loc}}(\Omega)\)，并且
\begin{equation}\label{b1:eq:sobolev-product}
 \partial_i(uv)=u\partial_i v+v\partial_i u
 \quad\text{几乎处处},\qquad 1\le i\le d.
\end{equation}
若两个函数均属于全域 \(W^{1,p}(\Omega)\cap L^\infty(\Omega)\)，则乘积也属于该空间，且
\[
 \|uv\|_{W^{1,p}}
 \le\|u\|_\infty\cdot\|v\|_{W^{1,p}}
      +\|v\|_\infty\cdot\|u\|_{W^{1,p}}.
\]
\end{proposition}
\begin{proof}
先设 \(p<\infty\)，固定 \(U,V\) 如前一引理。局部卷积给出光滑 \(u_j,v_j\)，它们在 \(U\) 中依 \(W^{1,p}\) 收敛，并可取一个子列，使两列函数几乎处处收敛。非负核保证 \(|u_j|\)、\(|v_j|\) 具有由原函数在 \(V\) 上的本性界给出的公共界。

因此 \(u_jv_j\to uv\) 于 \(L^p(U)\)。对一个导数项，分解
\[
 u_j\partial_i v_j-u\partial_i v
 =u_j(\partial_i v_j-\partial_i v)
   +(u_j-u)\partial_i v.
\]
第一项由统一有界性和导数的 \(L^p\) 收敛趋零；第二项由几乎处处收敛及以 \(C|\partial_i v|^p\) 为控制函数的控制收敛定理趋零。另一导数项同理。对经典乘积公式取极限，再遍历内部开集，就得到\eqref{b1:eq:sobolev-product}。

若 \(p=\infty\)，先把相同函数视为局部 \(W^{1,1}\) 成员，用刚证明的公式；其右端由假设局部本性有界，所以确实给出 \(W^{1,\infty}_{\mathrm{loc}}\) 关系。全域情形直接估计右端各分量，应用有限分量的 Minkowski 不等式，即得范数界；乘积的本性有界性也显然成立。
\end{proof}

只知道 \(u,v\in W^{1,p}\) 时，不能无条件断言乘积仍属于 \(W^{1,p}\)。弱导数公式并不自动提高 \(u\partial_i v\) 的可积性；第\ref{b1:ch:24}章的嵌入定理以后会提供另外一些充分条件。习题\ref{b1:ex:21-product-integrability-loss}则直接构造乘积失去原有可积性的例子。

若其中一个因子为光滑函数，不需要假设另一因子局部有界。第20章的弱 Leibniz 公式直接给出
\[
 \partial_i(au)=a\partial_i u+(\partial_i a)u.
\]
特别地，\(a\in C_c^\infty(\Omega)\) 时可以截出一个内部部分；这里产生的额外项 \(u\partial_i a\) 不能省略。

\subsection{链式法则}
\label{b1:subsec:210302}

本小节先讨论实值函数。对于向量值映射，把每个分量分别放在 Sobolev 空间中，证明完全相同；复值函数可以按两个实分量处理，而不是擅自要求复合映射全纯。

\begin{theorem}[链式法则]\label{b1:thm:sobolev-chain}
设 \(u\in W^{1,p}_{\mathrm{loc}}(\Omega;\mathbb R^m)\)，\(F\in C^1(\mathbb R^m;\mathbb R^\ell)\)，且 \(DF\) 有界。则
\[
 F\circ u\in W^{1,p}_{\mathrm{loc}}(\Omega;\mathbb R^\ell),
 \qquad
 \partial_i(F\circ u)=DF(u)\partial_i u
 \quad\text{几乎处处}.
\]
若 \(u\) 局部本性有界，则可去掉 \(DF\) 全域有界的要求。
\end{theorem}
\begin{proof}
有界导数使 \(F\) 为 Lipschitz 映射，因而
\[
 |F(u)|\le |F(0)|+\|DF\|_\infty\cdot|u|.
\]
右端局部属于 \(L^p\)，包括无穷端点。先取 \(p<\infty\)，并在固定内部区域用前面的卷积引理选取 \(u_j\to u\) 于 \(W^{1,p}\)，同时几乎处处收敛。Lipschitz 界给出 \(F(u_j)\to F(u)\) 于 \(L^p\)。导数分解为
\[
 DF(u_j)\partial_i u_j-DF(u)\partial_i u
 =DF(u_j)(\partial_i u_j-\partial_i u)
   +\bigl(DF(u_j)-DF(u)\bigr)\partial_i u.
\]
第一项的 \(L^p\) 范数趋零。第二项由于 \(DF\) 连续而几乎处处趋零，并受 \(2\|DF\|_\infty\cdot|\partial_i u|\) 控制，所以也依 \(L^p\) 趋零。经典链式法则在极限下给出弱导数公式。

\(p=\infty\) 时先在局部使用 \(p=1\) 的结论，再用所得公式的本性界。若 \(u\) 仅局部本性有界，对固定区域取包含其值域和卷积值域的紧球，用光滑截断把 \(F\) 改成该球邻域上与其相等、导数全域有界的映射，即归结为已证情形。
\end{proof}

全域结论还需要零阶项可积。若 \(u\in W^{1,p}(\Omega)\)，\(DF\) 有界且 \(F(0)=0\)，则 \(F(u)\in W^{1,p}(\Omega)\)。若 \(p<\infty\) 且 \(m_d(\Omega)<\infty\)，允许 \(F(0)\ne0\) 也可以；但在无限测度域上，\(u=0\)、\(F(t)=1+t\) 已经说明不能遗漏这个条件。\(p=\infty\) 时，全域本性有界的 \(u\) 经连续 \(F\) 映射后仍本性有界。

稍后要把光滑的 \(F\) 换成折线函数。折点处不能任意填写一个经典导数便结束论证，所需事实是：Sobolev 函数在一个水平集上没有非零弱梯度。

\begin{proposition}\label{b1:prop:sobolev-gradient-on-level-sets}
设实值 \(u\in W^{1,p}_{\mathrm{loc}}(\Omega)\)。对每个 \(c\in\mathbb R\)，有
\[
 \nabla u=0\quad\text{几乎处处于 }\{u=c\}.
\]
这里的水平集可以用任一可测代表定义，结论不受零测修改影响。
\end{proposition}
\begin{proof}
取 \(0\le\theta\le1\)、\(\theta\in C_c^\infty((-1,1))\)、\(\theta(0)=1\)，令
\[
 F_\varepsilon(t)=\int_c^t\theta\left(\frac{s-c}{\varepsilon}\right)\,\mathrm ds.
\]
有 \(|F_\varepsilon|\le2\varepsilon\)，而链式法则给出
\[
 \partial_i F_\varepsilon(u)
 =\theta\left(\frac{u-c}{\varepsilon}\right)\partial_i u.
\]
函数在每个内部有界区域趋于零，右端则由控制收敛在局部 \(L^1\) 中趋于 \(\mathbf1_{\{u=c\}}\partial_i u\)。弱导数关系的闭性说明后一个极限是零函数的弱导数，因此等于零。这个证明也包括 \(p=\infty\)，并没有使用尚待证明的沿直线代表或近似可微性。
\end{proof}

\subsection{截断}
\label{b1:subsec:210303}

对实数 \(M>0\)，定义截断函数
\[
 T_M(t)=\max\{-M,\min\{t,M\}\}.
\]
它保留 \([-M,M]\) 内的值，并把其余值压到两个端点。与把定义域截成一个子集不同，这是一种作用于函数值的操作。

\begin{proposition}\label{b1:prop:sobolev-truncation}
若实值 \(u\in W^{1,p}_{\mathrm{loc}}(\Omega)\)，则
\[
 T_M(u)\in W^{1,p}_{\mathrm{loc}}(\Omega),\qquad
 \nabla T_M(u)=\mathbf1_{\{|u|<M\}}\nabla u
 \quad\text{几乎处处}.
\]
正部 \(u^+=\max\{u,0\}\)、负部 \(u^-=\max\{-u,0\}\) 和绝对值也属于同一局部空间，且
\[
 \nabla u^+=\mathbf1_{\{u>0\}}\nabla u,\quad
 \nabla u^-=-\mathbf1_{\{u<0\}}\nabla u,\quad
 \nabla|u|=\operatorname{sgn}(u)\nabla u.
\]
约定 \(\operatorname{sgn}(0)=0\)。若 \(u\in W^{1,p}(\Omega)\)，上述函数也属于全域空间。
\end{proposition}
\begin{proof}
在实轴上用非负偶光滑核 \(\eta_\varepsilon\) 卷积 \(T_M\)，记为 \(F_\varepsilon\)。由于 \(T_M\) 奇且为 \(1\)-Lipschitz，\(F_\varepsilon(0)=0\)、\(|F_\varepsilon'|\le1\)，并且 \(F_\varepsilon\to T_M\) 一致成立。分段求导或直接积分可得
\[
 F_\varepsilon'=\mathbf1_{(-M,M)}*\eta_\varepsilon.
\]
因此当 \(t\ne\pm M\) 时，导数趋于 \(\mathbf1_{(-M,M)}(t)\)。在 \(\{u=\pm M\}\) 上，前一命题保证 \(\nabla u=0\)，所以无论近似导数在两个端点的极限为何，都有
\[
 F_\varepsilon'(u)\nabla u
 \longrightarrow\mathbf1_{\{|u|<M\}}\nabla u
 \quad\text{于 }L^1_{\mathrm{loc}}.
\]
由链式法则和弱导数的闭性得到所求关系；右端的局部 \(L^p\) 控制给出空间归属。

对 \(t^+\) 作同样卷积，并减去卷积在零点的值，使近似函数仍在零点为零。其导数在 \(t\ne0\) 处趋于 \(\mathbf1_{(0,\infty)}(t)\)，在零水平集上仍由梯度为零消除折点的不确定性。于是正部公式成立。负部及绝对值由 \((-u)^+\) 和 \(u^++u^-\) 得到。

最后，\(|T_M(u)|\)、\(u^+\)、\(u^-\) 和 \(|u|\) 均不超过 \(|u|\)，各导数的绝对值也由原导数控制，故全域情形不会丢失零阶或一阶可积性。
\end{proof}

若 \(1\le p<\infty\) 且 \(u\in W^{1,p}(\Omega)\)，则
\begin{equation}\label{b1:eq:sobolev-truncation-convergence}
 T_M(u)\longrightarrow u\quad\text{于 }W^{1,p}(\Omega)
 \quad(M\to\infty).
\end{equation}
事实上，函数误差由 \(|u|\mathbf1_{\{|u|>M\}}\) 控制，导数误差由 \(|\partial_i u|\mathbf1_{\{|u|\ge M\}}\) 控制，对它们的 \(p\) 次幂分别用控制收敛即可。在全域 \(W^{1,\infty}\) 中，对一个固定 \(u\)，只要 \(M\ge\|u\|_\infty\)，截断已经与原函数几乎处处相等；这不意味着正部等非线性运算对 \(W^{1,\infty}\) 范数总是连续，习题会区分这两个问题。

\subsection{坐标变换}
\label{b1:subsec:210304}

改变坐标还会改变积分域和零测集。若一个坐标映射把零测集拉回成正测度集，\(u\circ\Phi\) 就可能依赖 \(u\) 的代表；因此先明确映射的几何控制。

\begin{theorem}\label{b1:thm:sobolev-bilipschitz-change}
设 \(U,V\subset\mathbb R^d\) 为开集，\(\Phi:U\to V\) 是双 Lipschitz 同胚，且
\[
 |\Phi(x)-\Phi(y)|\le L|x-y|,\qquad
 |\Phi^{-1}(z)-\Phi^{-1}(w)|\le M|z-w|.
\]
若 \(u\in W^{1,p}(V)\)，则 \(u\circ\Phi\in W^{1,p}(U)\)，且
\begin{equation}\label{b1:eq:sobolev-coordinate-chain}
 \nabla(u\circ\Phi)(x)
   =D\Phi(x)^{\mathsf T}(\nabla u)(\Phi(x))
 \quad\text{对几乎处处的 }x\in U.
\end{equation}
包含零阶项的范数满足
\[
 \|u\circ\Phi\|_{W^{1,p}(U)}
 \le C(d,p,L,M)\|u\|_{W^{1,p}(V)}.
\]
\end{theorem}
\begin{proof}
由 Lipschitz 映射的零测集性质，\(\Phi^{-1}\) 把 \(V\) 中的 Lebesgue 零测集映成 \(U\) 中的零测集，所以复合不依赖代表。对非负可测 \(g\)，第17章的加权面积公式用于单射映射 \(\Phi^{-1}\)，给出
\begin{equation}\label{b1:eq:sobolev-change-measure-bound}
 \int_U g(\Phi(x))\,\mathrm dx
 =\int_V g(y)|\det D\Phi^{-1}(y)|\,\mathrm dy
 \le M^d\int_Vg(y)\,\mathrm dy.
\end{equation}
这里 Rademacher 定理保证导数几乎处处存在，算子范数不超过 \(M\)，从而行列式绝对值不超过 \(M^d\)。

先设 \(p<\infty\)，固定 \(U_0\) 满足 \(\overline{U_0}\Subset U\)。其像的闭包紧含于 \(V\)，因而可以在像的一个内部邻域对 \(u\) 作局部 \(W^{1,p}\) 光滑逼近 \(u_j\)。函数 \(u_j\circ\Phi\) 局部 Lipschitz；经典链式法则在 \(\Phi\) 的可微点给出
\[
 \nabla(u_j\circ\Phi)
 =D\Phi^{\mathsf T}(\nabla u_j)\circ\Phi.
\]
由\cref{b1:thm:lipschitz-weak-derivative}，这也是弱导数等式。估计\eqref{b1:eq:sobolev-change-measure-bound}使复合后的函数差及导数差在 \(U_0\) 中依 \(L^p\) 趋零；\(D\Phi\) 的算子范数又不超过 \(L\)，故右端也在 \(L^p\) 中收敛。弱导数关系的闭性于是给出\eqref{b1:eq:sobolev-coordinate-chain}。由于 \(U_0\) 任意，这一关系在整个 \(U\) 上成立。

若 \(p=\infty\)，同一局部证明先使用有限指数 \(1\)，得到公式后再取本性界。最后，对有限 \(p\) 用\eqref{b1:eq:sobolev-change-measure-bound}，对无穷端点用零测集拉回性质，并结合有限维向量范数等价性，得到包含函数与全部一阶偏导数的全域范数估计。
\end{proof}

对 \(\Phi^{-1}\) 再用一次定理可知，这个复合映射实际给出两个 \(W^{1,p}\) 空间之间的有界线性同构。它并不保持本书范数的数值，也不由单向 Lipschitz 控制推出逆方向的有界性。

高阶空间需要更强的坐标条件。若 \(\Phi\) 还是 \(C^k\) 微分同胚，并且 \(1\) 至 \(k\) 阶的各导数有界，则相同方法给出 \(u\circ\Phi\in W^{k,p}(U)\) 及相应有界估计：对光滑 \(u_j\)，逐次求导后，每一项都是某个 \((\partial^\beta u_j)\circ\Phi\) 乘以有限个 \(\Phi\) 的导数，且 \(|\beta|\le k\)；系数有统一界，故逐项用换元估计取极限即可。若还要得到反向的 \(W^{k,p}\) 有界映射，同样要控制 \(\Phi^{-1}\) 的高阶导数。仅有双 Lipschitz 性只控制一阶，不能代替这些条件。

% !TeX root = ../../Book1.tex
\section{Sobolev 函数的代表}
\label{b1:sec:2104}

Sobolev 空间按几乎处处相等取商，而连续性、沿直线的变化和差商都涉及点值。讨论这些性质时，必须先说明选择什么代表。本节先把一维积分表示推广到几乎所有坐标直线，再区分通常微分与近似微分。以下以实值函数为主；有限维向量值函数的相应结论可逐分量应用。

% 实读 Hajlasz2003SobolevMetric，作者公开正式版，Contemp. Math. 338 (2003)，
% §2 pp.175--177：ACL 定义、定理2.3及式(2.5)；§7定理7.13 pp.192--193；
% §9定理9.4及完整证明 pp.204--205。
% 原文7.13经上梯度并使用全域光滑稠密性；本节不移用该前置，
% ACL 改用第20章局部磨光及可数内矩形上的公共快速收敛子列。
% 双点估计改编9.4的均值望远镜路线，立方体平均振荡先由ACL自足证明。

\subsection{一维绝对连续代表}
\label{b1:subsec:210401}

一维情形中，弱导数不只决定积分配对，也决定代表在任意两点之间的增量。

\begin{proposition}\label{b1:prop:sobolev-interval-representative}
设 \(I\) 为开区间、\(1\leq p\leq\infty\)。函数 \(u\) 属于 \(W^{1,p}_{\mathrm{loc}}(I)\)，当且仅当它有局部绝对连续代表 \(\widetilde u\)，且 \(\widetilde u,\widetilde u'\in L^p_{\mathrm{loc}}(I)\)。这个代表唯一，经典导数与弱导数几乎处处相等。

若 \(I=(a,b)\) 有界且 \(u\in W^{1,p}(I)\)，则该代表还唯一延伸为 \([a,b]\) 上的绝对连续函数，并满足
\begin{equation}\label{b1:eq:interval-sobolev-sup}
 \|\widetilde u\|_{L^\infty([a,b])}
 \leq\dfrac{\|u\|_{L^1(I)}}{b-a}+\|u'\|_{L^1(I)}.
\end{equation}
\end{proposition}
\begin{proof}
局部刻画及唯一性就是定理\ref{b1:thm:one-dimensional-weak-ac}。在有界区间上，由 Hölder 不等式，\(u,u'\in L^1(I)\)。固定内点 \(x_0\)，积分表示
\[
 \widetilde u(x)=\widetilde u(x_0)+\int_{x_0}^x u'(t)\,\mathrm dt
\]
一直延伸到两端，并在整个闭区间上绝对连续。对任意 \(x,y\in[a,b]\)，有
\[
 |\widetilde u(x)|\leq|\widetilde u(y)|+\int_a^b|u'(t)|\,\mathrm dt.
\]
对 \(y\) 取平均，再对 \(x\) 取上确界，即得\eqref{b1:eq:interval-sobolev-sup}。
\end{proof}

当 \(1<p<\infty\) 时，同一积分表示给出
\[
 |\widetilde u(t)-\widetilde u(s)|
 \leq\|u'\|_{L^p(I)}\cdot|t-s|^{1-1/p}
 \quad(s,t\in[a,b]).
\]
当 \(p=\infty\) 时，右侧改为 \(\|u'\|_\infty\cdot|t-s|\)；当 \(p=1\) 时，对 \(s<t\) 的增量由 \(\int_s^t|u'|\) 控制，但不能从 \(L^1\) 范数单独读出一个正的 Hölder 指数。估计\eqref{b1:eq:interval-sobolev-sup}还有另一个用途：函数及其一阶导数的 \(L^1\) 收敛，会使对应的绝对连续代表一致收敛。下一小节将在几乎每条直线上使用这一点。

\subsection{ACL 性质}
\label{b1:subsec:210402}

在多维开集内，平行于某个坐标轴的直线可能与区域相交为多个区间，所以这里的绝对连续性应逐个紧子区间理解。先设 \(d\geq2\)；\(d=1\) 时，下面的条件就是局部绝对连续。

\begin{definition}\label{b1:def:acl}
设 \(v:\Omega\rightarrow\mathbb R\) 为有限值 Borel 函数。若对每个坐标方向 \(e_i\)，对其正交补中几乎所有 \(z\)，限制
\[
 t\longmapsto v(z+te_i)
\]
在开集 \(\{\,t\in\mathbb R\mid z+te_i\in\Omega\,\}\) 的每个紧子区间上绝对连续，则称 \(v\) \term[yan-zuobiaoxian-jueduilianxu]{沿坐标线绝对连续}[absolutely continuous on lines]，简称具有 ACL 性质。这里的“几乎所有”使用横向变量 \(z\) 上的 \((d-1)\) 维 Lebesgue 测度。
\end{definition}

\begin{theorem}\label{b1:thm:sobolev-acl}
设 \(1\leq p\leq\infty\)，且 \(u,g_1,\ldots,g_d\in L^p_{\mathrm{loc}}(\Omega)\)。则 \(\partial_i u=g_i\) 对所有 \(i\) 成立，当且仅当 \(u\) 有一个具有 ACL 性质的代表 \(\widetilde u\)，其经典坐标偏导数几乎处处存在，且等于 \(g_i\)。
\end{theorem}
\begin{proof}
先设弱导数关系成立。列出所有闭包紧含于 \(\Omega\) 的有理端点开矩形 \(Q_m\)。由\cref{b1:lem:sobolev-local-smoothing}，可以选 \(\varepsilon_j\downarrow0\)，使 \(u_j=u*\rho_{\varepsilon_j}\) 在 \(Q_1,\ldots,Q_j\) 的邻域上都有定义，且
\[
 \|u_j-u\|_{L^1(Q_m)}
 +\sum_{i=1}^d\|\partial_i u_j-g_i\|_{L^1(Q_m)}<2^{-j}
 \quad(m\leq j).
\]
这里仅用了内矩形上的局部 \(L^1\) 逼近。将每个 \(u_j\) 在其定义域外补零，只为得到 \(\Omega\) 上的 Borel 函数；不宣称补零后的函数在全域光滑或保留弱导数。

固定 \(Q_m\) 和方向 \(i\)，写 \(Q_m=Q_m'\times I\)，把第 \(i\) 个坐标放在最后。对上述可求和误差应用 Tonelli 定理，可知对几乎每个 \(z\in Q_m'\)，
\[
 \sum_{j\geq m}\left(
  \|u_j(z,\cdot)-u(z,\cdot)\|_{L^1(I)}
  +\|\partial_i u_j(z,\cdot)-g_i(z,\cdot)\|_{L^1(I)}
 \right)<\infty.
\]
因此这条直线上的 \(u_j\) 在一维 \(W^{1,1}(I)\) 中为 Cauchy 列。将\eqref{b1:eq:interval-sobolev-sup}用于差 \(u_j-u_k\)，得到它们在 \(I\) 上一致收敛到某个函数 \(v_z\)。又因导数在 \(L^1(I)\) 中趋于 \(g_i(z,\cdot)\)，在光滑函数的积分公式中取极限，得
\[
 v_z(t)-v_z(s)=\int_s^t g_i(z,r)\,\mathrm dr
 \quad(s,t\in I).
\]
所以 \(v_z\) 绝对连续，导数几乎处处为该截面的 \(g_i\)。它也与 \(u(z,\cdot)\) 几乎处处相等。

现在在 \(\Omega\) 中统一定义
\[
 \widetilde u(x)=
 \begin{cases}
  \displaystyle\lim_{j\to\infty}u_j(x),&\text{若极限存在且有限},\\
  0,&\text{其他情形}.
 \end{cases}
\]
这是一个 Borel 函数，并由各矩形上的可求和 \(L^1\) 误差知 \(\widetilde u=u\) 几乎处处。在前述每条良好直线上，它在矩形内的每一点都等于一致极限 \(v_z\)。对可数个矩形及有限个方向合并横向零测例外后，仍只排除了各方向中的一个横向零测集。任意紧线段可用有限个这样的矩形覆盖，局部积分表示在交叠处来自同一个函数 \(\widetilde u\)，故可拼接。这证明了同一个代表同时具有全部坐标方向的 ACL 性质。

反过来，在内矩形中取测试函数，沿第 \(i\) 个方向在几乎每条良好直线上作一维分部积分，再由 Fubini 定理得
\[
 \int_{Q_m}\widetilde u\,\partial_i\varphi
 =-\int_{Q_m}g_i\varphi.
\]
各积分因局部可积性而有定义；测试函数在相应区间端点附近为零。弱导数的局部性随后把这些关系拼成整个 \(\Omega\) 上的关系。
\end{proof}

沿几乎所有直线成立，与沿每一条直线成立不同。维数至少为二时，甚至把一个 ACL 代表在单点改值，也只会破坏有限条坐标直线上的绝对连续性，而不改变 ACL 性质。因此 ACL 代表不像一维连续代表那样逐点唯一。定理所需的是存在一个共同代表，并不是给等价类的每一点都指定不可更改的值。

\subsection{\texorpdfstring{\(W^{1,\infty}\)}{W1,infinity} 与局部 Lipschitz 代表}
\label{b1:subsec:210403}

由 ACL 性质，Sobolev 函数的合适代表有几乎处处存在的坐标偏导数，且它们组成弱梯度。但坐标偏导存在不等于通常的全微分存在。本小节先给出能够恢复通常微分的有界梯度情形。

\begin{theorem}\label{b1:thm:w1infty-lipschitz}
函数 \(u\) 属于 \(W^{1,\infty}_{\mathrm{loc}}(\Omega)\)，当且仅当它有局部 Lipschitz 代表。该连续代表唯一，并几乎处处通常可微，其经典微分为弱梯度所确定的线性映射。

更具体地，若 \(B\) 为开球、\(\overline B\Subset\Omega\)，且 \(|\nabla u|\leq L\) 几乎处处于 \(B\)，则此代表在 \(B\) 上为 \(L\)-Lipschitz 函数。
\end{theorem}

这里 \(|\nabla u|\) 使用 \(\mathbb R^d\) 的欧氏范数。它与本章按偏导分量取最大值的范数满足
\[
 \bigl\||\nabla u|\bigr\|_{L^\infty(B)}
 \leq \sqrt d\,\max_{1\leq i\leq d}\|\partial_i u\|_{L^\infty(B)}
 \leq \sqrt d\,\|u\|_{W^{1,\infty}(B)}.
\]
所以若只从本章的 \(W^{1,\infty}\) 范数读取梯度控制，可以直接得到 \(\sqrt d\) 倍的 Lipschitz 常数；定理中的 \(L\) 则是欧氏梯度模的本性上界，不应与前一常数混同。

\begin{proof}
先设 \(u\in W^{1,\infty}_{\mathrm{loc}}\)。固定所述开球，令 \(u_\varepsilon=u*\rho_\varepsilon\)。局部卷积微分公式给出
\[
 \nabla u_\varepsilon=(\nabla u)*\rho_\varepsilon,
 \quad |\nabla u_\varepsilon|\leq L
\]
于卷积只涉及 \(B\) 内点的区域。取 \(B\) 中两个 Lebesgue 点 \(x,y\)。线段 \([x,y]\) 紧含于 \(B\)，故当 \(\varepsilon\) 足够小时，在这条线段上积分光滑函数的导数可得
\[
 |u_\varepsilon(x)-u_\varepsilon(y)|\leq L|x-y|.
\]
由\cref{b1:prop:approx-identity-lebesgue-point}，两端趋于 \(u\) 在相应点的 Lebesgue 值。于是同一 Lipschitz 不等式在 \(B\) 的 Lebesgue 点集上成立。

该点集满测度，因而在 \(B\) 中稠密。对任意点取其中的逼近点列，Lipschitz 估计使函数值为 Cauchy 列；估计也保证极限不依赖所选点列，并在整个 \(B\) 上保留常数 \(L\)。所得连续函数与 \(u\) 几乎处处相等。在交叠球上，两个连续代表几乎处处相等，所以处处相等，由此拼成局部 Lipschitz 代表。唯一性同样来自连续性与几乎处处相等。

反向由定理\ref{b1:thm:lipschitz-weak-derivative}得到。最后对局部 Lipschitz 代表使用 Rademacher 定理\ref{b1:thm:rademacher}；经典偏导与弱偏导几乎处处一致，故全微分的系数正是弱梯度。
\end{proof}

这里的“局部”不能无条件去掉。例如在 \(\Omega=(-1,0)\cup(0,1)\) 上，令 \(u\) 左侧为零、右侧为一，则 \(u\in W^{1,\infty}(\Omega)\)，弱导数为零；但跨越缺失的原点取两侧点，差商可以任意大，故它不是 \(\Omega\) 上的 Lipschitz 函数。

即使函数属于 \(W^{1,\infty}\)，通常可微的断言也针对所选的局部 Lipschitz 代表，而非任意代表。在有界开集上，\(\mathbf1_{\mathbb Q^d}\) 代表零元素，却处处不连续。一般有限 \(p\) 时，上述有界梯度证明不适用；不能把本小节标题理解成“所有 Sobolev 代表都几乎处处通常可微”。对全部 \(1\leq p\leq\infty\) 统一成立的结论是下一小节的近似可微性。

\subsection{近似可微性}
\label{b1:subsec:210404}

近似微分允许忽略密度为零的异常点，但仍要求一阶误差相对于距离趋于零。本小节的证明链为
\[
 \begin{gathered}
  \text{ACL}\Longrightarrow\text{立方体平均振荡估计}\\
  \Longrightarrow\text{极大函数双点估计}\\
  \Longrightarrow\text{可数 Lipschitz 分片}\\
  \Longrightarrow\operatorname{ap}Du=\nabla u.
 \end{gathered}
\]
其中层集 \(\{Mg\leq n\}\) 的作用是在每一层上得到 Lipschitz 分片，并不是把原函数提升为整个邻域上的 Lipschitz 函数。为从弱导数得到上述控制，先把沿坐标线的积分公式转成局部平均振荡的估计。对有限正体积的集合 \(Q\)，简记
\[
 u_Q=\dfrac1{m_d(Q)}\int_Q u.
\]

\begin{lemma}\label{b1:lem:sobolev-cube-oscillation}
设 \(u\in W^{1,1}_{\mathrm{loc}}(\Omega)\)，\(Q\) 为边长 \(\ell\) 的坐标开立方体，且 \(\overline Q\Subset\Omega\)。则
\begin{equation}\label{b1:eq:sobolev-cube-oscillation}
 \int_Q|u-u_Q|
 \leq\ell\sum_{i=1}^d\int_Q|\partial_i u|.
\end{equation}
\end{lemma}
\begin{proof}
选取 ACL 代表。对几乎所有 \(x,y\in Q\)，依次把 \(x\) 的各坐标换成 \(y\) 的坐标，用 \(d\) 条坐标线段连接它们。所有这些线段都属于良好直线的例外集之外；此断言对几乎所有点对成立，正是对每个方向的横向零测集应用 Fubini 定理的结果。沿线段积分，再用三角不等式，得到
\[
 |u(x)-u(y)|
 \leq\sum_{i=1}^d\int_{[x_i,y_i]}
 |\partial_i u(y_1,\ldots,y_{i-1},t,x_{i+1},\ldots,x_d)|\,\mathrm dt,
\]
其中 \([x_i,y_i]\) 表示两端点之间的无向区间。将此式对 \(x,y\) 积分。第 \(i\) 项中未出现的 \(d-1\) 个坐标积分产生因子 \(\ell^{d-1}\)；对固定 \(t\)，使区间 \([x_i,y_i]\) 包含 \(t\) 的点对面积至多为 \(\ell^2\)。由 Tonelli 定理，第 \(i\) 项的二重积分不超过 \(\ell^{d+1}\int_Q|\partial_i u|\)。因此
\[
 \int_Q|u(x)-u_Q|\,\mathrm dx
 \leq\dfrac1{m_d(Q)}\int_Q\int_Q|u(x)-u(y)|\,\mathrm dy\,\mathrm dx
 \leq\ell\sum_{i=1}^d\int_Q|\partial_i u|,
\]
因为 \(m_d(Q)=\ell^d\)。
\end{proof}

下一步比较两个点附近的一列局部均值。Hajłasz 的《Sobolev spaces on metric-measure spaces》\cite[Section 9, Theorem 9.4]{Hajlasz2003SobolevMetric}用这种逐半径的望远镜求和，从平均振荡得到双点控制；以下在欧氏立方体上给出所需证明。

\begin{proposition}\label{b1:prop:sobolev-maximal-difference}
设 \(w\in W^{1,1}(\mathbb R^d)\)，令 \(g=\sum_{i=1}^d|\partial_iw|\)。存在只依赖于 \(d\) 的常数 \(C_d\)，使在 \(w\) 的任意两个 Lebesgue 点 \(x,y\) 上，其 Lebesgue 值满足
\begin{equation}\label{b1:eq:sobolev-maximal-difference}
 |w(x)-w(y)|\leq C_d|x-y|\bigl(Mg(x)+Mg(y)\bigr).
\end{equation}
这里 \(M\) 是\secref{b1:sec:1303}的中心极大算子。
\end{proposition}
\begin{proof}
只需考虑 \(x\ne y\) 且两处极大函数都有限的情形。记 \(r=|x-y|\)，\(Q(x,t)=x+(-t,t)^d\)，并置 \(Q_j=Q(x,2^{-j}r)\)。由于立方体包含在可比半径的球内，Lebesgue 点性质给出 \(w_{Q_j}\rightarrow w(x)\)。由\eqref{b1:eq:sobolev-cube-oscillation}，
\[
 \begin{aligned}
 |w_{Q_{j+1}}-w_{Q_j}|
 &\leq\dfrac1{m_d(Q_{j+1})}\int_{Q_j}|w-w_{Q_j}|\\
 &\leq C_d\,2^{-j}r\,\dfrac1{m_d(Q_j)}\int_{Q_j}g
 \leq C_d\,2^{-j}r\,Mg(x).
 \end{aligned}
\]
最后一步使用 \(Q(x,t)\subseteq B(x,\sqrt d\,t)\) 及两者体积比只依赖维数。求和得到
\[
 |w(x)-w_{Q(x,r)}|\leq C_drMg(x).
\]
在 \(y\) 处也有同样的估计。

还需比较半径为 \(r\) 的两个均值。因为 \(Q(x,r)\) 和 \(Q(y,r)\) 都包含在 \(R=Q(x,2r)\) 中，体积比有统一界，故
\[
 \begin{aligned}
 |w_{Q(x,r)}-w_{Q(y,r)}|
 &\leq|w_{Q(x,r)}-w_R|+|w_{Q(y,r)}-w_R|\\
 &\leq\dfrac{C_d}{m_d(R)}\int_R|w-w_R|
 \leq C_drMg(x).
 \end{aligned}
\]
合并三项便得到\eqref{b1:eq:sobolev-maximal-difference}。
\end{proof}

极大函数在某个层集上有界时，上式就成为该层集上的 Lipschitz 估计。这里不需要 \(M\) 在 \(L^1\) 中强有界，只需要它对可积函数几乎处处有限，因此端点 \(p=1\) 也能纳入论证。

\begin{theorem}\label{b1:thm:sobolev-approximate-differentiability}
设 \(1\leq p\leq\infty\)，且 \(u\in W^{1,p}_{\mathrm{loc}}(\Omega)\)。则任一有限值可测代表都几乎处处近似可微，并且
\[
 \operatorname{ap}Du(x)h=\nabla u(x)\cdot h
 \quad(h\in\mathbb R^d)
\]
在几乎所有 \(x\in\Omega\) 上成立。
\end{theorem}

这里“任一代表”的量词是：每选定一个有限值可测代表，都存在一个可以依赖该代表的零测例外集，使结论在其外成立。不同代表只在零测集上不同，因而“几乎处处近似可微”作为等价类性质保持不变；这不表示所有代表在同一个共同点集上都近似可微。

\begin{proof}
局部 \(L^p\) 包含于局部 \(L^1\)，所以先把 \(u\) 看成 \(W^{1,1}_{\mathrm{loc}}\) 函数。固定 \(K\Subset\Omega\)，选 \(\chi\in C_c^\infty(\Omega)\)，使 \(\chi=1\) 于 \(K\) 的某个邻域。由光滑乘法公式，将 \(w=\chi u\) 及其弱导数补零到 \(\mathbb R^d\)，得到 \(w\in W^{1,1}(\mathbb R^d)\)。这一步的补零没有边界项，因为 \(\chi\) 紧支撑于 \(\Omega\)。

令 \(g=\sum_i|\partial_iw|\in L^1(\mathbb R^d)\)。由 Hardy--Littlewood 弱型不等式\ref{b1:thm:hl-weak-11}，\(Mg\) 几乎处处有限。从 \(K\) 删去零测例外，使剩下的点都是 \(w\) 的 Lebesgue 点，且所选的 \(u\) 值等于该点的 Lebesgue 值。再按 \(Mg(x)\leq n\) 分层，得可测集 \(E_n\)，其并覆盖 \(K\) 中几乎所有点。由\eqref{b1:eq:sobolev-maximal-difference}，\(u|_{E_n}\) 为 \(2C_dn\)-Lipschitz 函数。

每个 \(E_n\) 体积有限。由紧内正则性，可在其中选择可数个紧集，除去一个零测集后覆盖 \(E_n\)。再用可数个 \(K\Subset\Omega\) 覆盖 \(\Omega\)，就得到几乎处处覆盖定义域的紧集族 \((K_j)\)，每个 \(u|_{K_j}\) 都是 Lipschitz 的。由\cref{b1:thm:approximate-lipschitz-pieces}，\(u\) 几乎处处近似可微。

还要识别近似微分，而不只是证明它存在。把 \(u|_{K_j}\) 延拓为全空间 Lipschitz 函数 \(F_j\)。由定理\ref{b1:thm:lipschitz-weak-derivative}，\(F_j\) 的弱梯度等于其几乎处处经典梯度。又因 \(u-F_j=0\) 于 \(K_j\)，命题\ref{b1:prop:sobolev-gradient-on-level-sets}给出
\[
 \nabla u=\nabla F_j=DF_j
 \quad\text{几乎处处于 }K_j.
\]
其中最后一项按行向量与线性泛函的对应理解。在 \(K_j\) 的密度点且 \(F_j\) 可微的点上，沿 \(K_j\) 两函数相等，所以密度一集合判据给出 \(\operatorname{ap}Du=DF_j\)。合并这些等式，再合并可数个零测例外，便得所需识别。
\end{proof}

这个结论把弱导数与点态的一阶几何联系起来，但所得 Lipschitz 常数依赖于所选片，不能据此推出整个邻域上的 Lipschitz 性质。它也不声称 Sobolev 映射把每个零测集映成零测集：近似微分忽略定义域中的零测改值，像的测度却可能受到这些改值影响。以后将微分性质用于面积、迹或嵌入问题时，仍须分别检查那些问题需要的代表与额外条件。

% 本章小结（不计入编号）
\begin{chaptersummary}
Sobolev 空间的基本结构来自两个事实：各阶导数可以作为有限个 \(L^p\) 分量，而弱导数关系在这些分量的适当极限下保持。因此函数与导数不能独立指定，却可以一起使用 Banach 空间和弱收敛的方法。

\ProseParagraphBreak

运算规则首先是局部的。光滑逼近让经典公式进入弱导数框架，可积性条件再决定结果属于哪个全域空间。水平集上的梯度为零消除了截断函数折点处的不确定性；双 Lipschitz 坐标变换既控制导数，也保证零测集与积分估计正确转移。

\ProseParagraphBreak

对函数代表，一维绝对连续性、多维沿坐标线绝对连续性和局部 Lipschitz 性是不同层次。一般 Sobolev 函数几乎处处的近似微分等于弱梯度，但这不提供无条件的全域连续代表或边界值。下一章将处理全域光滑逼近、零边界空间与延拓，为以后讨论迹、嵌入和紧性准备工具。
\end{chaptersummary}


\BookExercises

% \chapref{b1:ch:21}习题临时稿；由章聚合文件已有的 BookExercises 入口统一接入。
% 十一题均为依据已读本章纲要、实际前章正文与已确认正文接口自拟的工具变式。
% 不冒认教材原题编号；不用第22章全局稠密性、第24章嵌入或尚未建立的迹理论。

\begin{exercise}\label{b1:ex:21-basic-definition-norm}
令 \(Q=(0,1)^2\)，\(u(x,y)=xy\)。对每个 \(1\leq p\leq\infty\)，直接从弱导数定义验证 \(u\in W^{1,p}(Q)\)，求出两个弱偏导数，并计算本章约定的 \(\|u\|_{W^{1,p}(Q)}\)。
\end{exercise}
% 来源：自拟；用于训练弱导数定义与本章逐分量范数，不调用后续嵌入或正则性。
\begin{solution}
取任意 \(\varphi\in C_c^\infty(Q)\)。由 Fubini 定理和一维分部积分，
\[
 \int_Qxy\,\partial_x\varphi(x,y)\,\mathrm dx\,\mathrm dy
 =\int_0^1y\left(\int_0^1x\,\partial_x\varphi(x,y)\,\mathrm dx\right)\mathrm dy
 =-\int_Qy\varphi(x,y)\,\mathrm dx\,\mathrm dy.
\]
测试函数在端点附近为零，所以没有边界项。交换两个变量得到
\[
 \partial_xu=y,\qquad \partial_yu=x
\]
的弱导数等式。三者在有界方形上都属于每个 \(L^p\)，故 \(u\in W^{1,p}(Q)\)。

若 \(1\leq p<\infty\)，则
\[
 \|u\|_p^p=\frac1{(p+1)^2},\qquad
 \|\partial_xu\|_p^p=\|\partial_yu\|_p^p=\frac1{p+1}.
\]
因此
\[
 \|u\|_{W^{1,p}(Q)}
 =\left(\frac1{(p+1)^2}+\frac2{p+1}\right)^{1/p}.
\]
当 \(p=\infty\) 时，\(u\)、\(x\) 与 \(y\) 的本性上确界都为 \(1\)，而本章范数取三者的最大值，故 \(\|u\|_{W^{1,\infty}(Q)}=1\)。
\end{solution}

\begin{exercise}\label{b1:ex:21-basic-vector-chain}
设 \(\Omega\subset\mathbb R^d\) 为开集、\(1\leq p\leq\infty\)，且
\[
 u=(u^1,u^2)\in W^{1,p}_{\mathrm{loc}}(\Omega;\mathbb R^2).
\]
令 \(F(s,t)=\sin s+\cos t\)。证明 \(F\circ u\in W^{1,p}_{\mathrm{loc}}(\Omega)\)，并对每个 \(1\leq i\leq d\) 求出 \(\partial_i(F\circ u)\)。
\end{exercise}
% 来源：自拟；直接训练向量值记号及正文链式法则的基础计算。
\begin{solution}
映射 \(F\in C^1(\mathbb R^2)\)，且
\[
 DF(s,t)=(\cos s,-\sin t)
\]
一致有界。由\cref{b1:thm:sobolev-chain}，复合函数属于所述局部 Sobolev 空间，并且
\[
 \partial_i(F\circ u)
 =\cos(u^1)\,\partial_i u^1
  -\sin(u^2)\,\partial_i u^2
 \quad\text{几乎处处于 }\Omega.
\]
右端属于 \(L^p_{\mathrm{loc}}\)，因为两个三角函数因子的绝对值不超过 \(1\)。局部陈述不要求 \(F\circ u\) 在无界域上具有全域 \(L^p\) 范数。
\end{solution}

\begin{exercise}\label{b1:ex:21-basic-derivative-embedding}
在 \(I=(0,1)\) 上令
\[
 u_n(x)=x+n^{-2}\sin(nx),\qquad u(x)=x.
\]
证明对每个 \(1\leq p\leq\infty\)，都有 \(u_n\to u\) 于 \(W^{1,p}(I)\)。再用\cref{b1:prop:sobolev-derivative-embedding}中的映射 \(\mathcal J\)，解释为什么这里的 Sobolev 收敛恰好等价于函数分量与一阶导数分量同时在 \(L^p(I)\) 中收敛。
\end{exercise}
% 来源：自拟；用于训练有限乘积嵌入和Sobolev范数收敛的直接判定。
\begin{solution}
两函数光滑，且
\[
 u_n-u=n^{-2}\sin(nx),\qquad
 u_n'-u'=n^{-1}\cos(nx).
\]
由于 \(m_1(I)=1\)，对每个有限 \(p\) 有
\[
 \|u_n-u\|_p\leq n^{-2},\qquad
 \|u_n'-u'\|_p\leq n^{-1},
\]
所以
\[
 \|u_n-u\|_{W^{1,p}(I)}
 \leq\bigl(n^{-2p}+n^{-p}\bigr)^{1/p}\longrightarrow0.
\]
在 \(p=\infty\) 时，本章范数取两个分量范数的最大值，同样有
\[
 \|u_n-u\|_{W^{1,\infty}(I)}\leq n^{-1}\longrightarrow0.
\]
这里 \(\mathcal Jv=(v,v')\)，而 \(Y_p\) 的乘积范数正按有限 \(p\) 的 \(p\) 次和或无穷端点的最大值定义。因此 \(\|\mathcal J(u_n-u)\|_{Y_p}=\|u_n-u\|_{W^{1,p}}\)，所述两个分量同时收敛与 Sobolev 范数收敛完全等价。
\end{solution}

\begin{optionalexercise}\label{b1:ex:21-scaling-concentration}
设整数 \(d\ge2\)、\(1\le p<d\)，记 \(B=B(0,1)\)，取非零 \(\psi\in C_c^\infty(B)\)，并令
\[
 u_r(x)=r^{1-d/p}\psi(x/r),\quad 0<r<1.
\]
计算任意多指标 \(\alpha\) 对应的 \(\|\partial^\alpha u_r\|_{L^p(B)}\)。证明 \(u_r\to0\) 于 \(L^p(B)\)，其一阶导数范数保持不变，而 \(\|u_r\|_\infty\to\infty\)。据此证明，不存在独立于 \(u\) 的常数 \(C\)，使
\[
 \|u\|_{L^\infty(B)}
 \le C\|u\|_{W^{1,p}(B)}
 \quad\bigl(u\in C_c^\infty(B)\bigr).
\]
\end{optionalexercise}
% 来源：自拟；由弱导数与线性换元的尺度计算构造集中列，只排除p<d的所列估计，不预用嵌入定理。
\begin{solution}
函数 \(u_r\) 的支集包含在 \(B(0,r)\) 中，经典求导与换元 \(x=ry\) 给出
\[
 \begin{aligned}
 \|\partial^\alpha u_r\|_p^p
 &=r^{p-d-p|\alpha|}\int_{B(0,r)}
       |(\partial^\alpha\psi)(x/r)|^p\,\mathrm dx\\
 &=r^{p(1-|\alpha|)}\|\partial^\alpha\psi\|_p^p.
 \end{aligned}
\]
因此
\[
 \|\partial^\alpha u_r\|_p
 =r^{1-|\alpha|}\|\partial^\alpha\psi\|_p.
\]
当 \(|\alpha|=0\) 时右端趋零；当 \(|\alpha|=1\) 时它不随 \(r\) 改变。由于非零紧支集光滑函数不可能在连通的全空间上为常数，至少一个一阶导数的范数为正。故这列函数的 \(W^{1,p}\) 范数一致有界，而一阶导数并没有一同趋零。

另一方面，\(\psi\) 连续，故其上确界等于本性上确界，且
\[
 \|u_r\|_\infty=r^{1-d/p}\|\psi\|_\infty\longrightarrow\infty,
\]
因为 \(p<d\)。若题中统一估计成立，左端便会被一致有界的 Sobolev 范数控制，矛盾。这个例子只处理严格不等式 \(p<d\)；当 \(p=d\) 时，同一伸缩的振幅不再增长，不能从本题推出临界指数的结论。
\end{solution}

\begin{optionalexercise}\label{b1:ex:21-endpoint-structure}
\begin{enumerate}
\item\label{b1:item:21-w11-no-weak-subsequence}
在 \(I=(0,1)\) 上令 \(v_n(x)=\min(nx,1)\)。证明 \((v_n)\) 在 \(W^{1,1}(I)\) 中有界，但没有弱收敛子列。
\item\label{b1:item:21-w1infty-nonseparable}
对 \(t\in(0,1)\)，令 \(w_t(x)=|x-t|\)。证明这些元素在 \(W^{1,\infty}(I)\) 中两两距离至少为 \(2\)，并据此证明该空间不可分。
\end{enumerate}
\end{optionalexercise}
% 来源：自拟；第10章端点弱测试与本章连续导数映射的具体应用，完整折点构造未在正文展开。
\begin{solution}
第一问中，\(v_n\) 为 Lipschitz 函数，弱导数几乎处处等于经典导数，所以
\[
 v_n'=n\mathbf1_{(0,1/n)},\quad
 \|v_n\|_1\le1,\quad \|v_n'\|_1=1.
\]
这证明了 Sobolev 范数有界。

假设某个子列 \((v_{n_j})\) 在 \(W^{1,1}(I)\) 中弱收敛。由\cref{b1:prop:sobolev-derivative-embedding}，导数映射从 \(W^{1,1}(I)\) 到 \(L^1(I)\) 有界线性，相应导数必在 \(L^1(I)\) 中弱收敛到某个 \(g\)。固定 \(a\in(0,1)\)，以及任意在 \((0,a)\) 上几乎处处为零的 \(h\in L^\infty(I)\)。当 \(n_j>1/a\) 时，\(\int_Iv_{n_j}'h=0\)，故 \(\int_Igh=0\)。这些有界测试说明 \(g=0\) 几乎处处于 \((a,1)\)。取 \(a=1/k\)、\(k\ge2\)，便得 \(g=0\) 几乎处处于 \(I\)。但常函数 \(h=1\) 也是允许的对偶测试，且
\[
 \int_Ig=\lim_j\int_Iv_{n_j}'=1,
\]
矛盾。因此不存在所假设的弱收敛子列。事实上，\(v_n\to1\) 于 \(L^1(I)\)，失去弱紧性的是向端点集中的导数。

第二问中，\(w_t'=\operatorname{sgn}(x-t)\) 几乎处处。若 \(s<t\)，则在正测度区间 \((s,t)\) 上有
\[
 w_s'=1,\quad w_t'=-1.
\]
所以 \(\|w_s'-w_t'\|_\infty=2\)，从而 \(\|w_s-w_t\|_{W^{1,\infty}}\ge2\)。

若空间有可数稠密子集 \(\{z_j\}\)，则每个开球 \(B(w_t,1/2)\) 至少包含一个 \(z_j\)。这些球两两不交，否则三角不等式将使两个中心的距离小于 \(1\)。这样便给不可数多个球分别配上不同的可数集元素，矛盾。故 \(W^{1,\infty}(I)\) 不可分。
\end{solution}

\begin{exercise}\label{b1:ex:21-oscillatory-gradient-energy}
在 \(I=(0,2\pi)\) 上令 \(u_n(x)=\sin(nx)/n\)。证明 \(u_n\rightharpoonup0\) 于 \(H^1(I)\)，但不依 \(H^1\) 范数趋零。进一步证明，对任意 \(g\in C([0,2\pi])\)，
\[
 \lim_{n\to\infty}\int_0^{2\pi}g(x)|u_n'(x)|^2\,\mathrm dx
 =\frac12\int_0^{2\pi}g(x)\,\mathrm dx.
\]
解答不预借 Fourier 理论。
\end{exercise}
% 来源：自拟；与第20章分布收敛例相比，新增H1弱极限识别及导数能量的非零极限。
\begin{solution}
先证明 \(\cos(nx)\rightharpoonup0\) 于 \(L^2(I)\)。对一个子区间 \((a,b)\subset I\)，
\[
 \left|\int_a^b\cos(nx)\,\mathrm dx\right|
 =\frac{|\sin(nb)-\sin(na)|}{n}\le\frac2n.
\]
因此与任意有限区间阶梯函数相乘后的积分都趋于零。这些阶梯函数在 \(L^2(I)\) 中稠密：先用\secref{b1:sec:1002}的连续紧支集函数逼近，再在有限区间上利用一致连续性作阶梯逼近。又因 \(\|\cos(nx)\|_2=\sqrt\pi\)，对 \(h\in L^2(I)\) 及一个阶梯函数 \(s\)，有
\[
 \left|\int_I\cos(nx)h\right|
 \le\left|\int_I\cos(nx)s\right|
     +\sqrt\pi\|h-s\|_2.
\]
先选 \(s\)，再令 \(n\to\infty\)，即得所需弱收敛。

现在 \(\|u_n\|_2=\sqrt\pi/n\to0\)，而 \(u_n'=\cos(nx)\rightharpoonup0\) 于 \(L^2\)。由\cref{b1:prop:sobolev-weak-components}，\(u_n\rightharpoonup0\) 于 \(H^1\)。也可以直接对任意 \(v\in H^1(I)\) 测试内积：零阶积分由 \(\|u_n\|_2\to0\) 控制，一阶积分由刚证的弱收敛控制。不过
\[
 \|u_n'\|_2^2=\pi
\]
始终不变，所以不可能强收敛到零。

最后使用恒等式 \(\cos^2(nx)=\bigl(1+\cos(2nx)\bigr)/2\)。函数 \(g\) 属于 \(L^2(I)\)，上面的弱收敛结论同样适用于 \(\cos(2nx)\)，于是
\[
 \int_Ig|u_n'|^2
 =\frac12\int_Ig+\frac12\int_Ig\cos(2nx)
 \longrightarrow\frac12\int_Ig.
\]
函数值的强收敛与导数的弱收敛都不能使导数平方直接趋向极限导数的平方。
\end{solution}

\begin{exercise}\label{b1:ex:21-product-integrability-loss}
在 \(B=B(0,1)\subset\mathbb R^3\) 上取 \(u(x)=|x|^{-1/3}\)，原点处任意赋一个有限值。证明
\[
 u\in H^1(B),\quad
 u^2\in L^2(B)\cap W^{1,1}(B),\quad
 u^2\notin H^1(B).
\]
说明这与有界 Sobolev 函数的乘积法则并不矛盾。
\end{exercise}
% 来源：自拟；根正文只说明W1p一般非代数，本题独立计算高维奇性，用ACL反向判据确认弱导数。
\begin{solution}
在原点外，经典梯度为
\[
 \nabla u(x)=-\frac13|x|^{-7/3}x,\quad
 \nabla(u^2)(x)=-\frac23|x|^{-8/3}x.
\]
所以其大小分别为 \(\tfrac13|x|^{-4/3}\) 与 \(\tfrac23|x|^{-5/3}\)。由三维径向积分公式，\(|x|^{-q}\) 在单位球上可积，当且仅当 \(q<3\)，因为相应径向积分是常数 \(4\pi\) 乘以 \(\int_0^1r^{2-q}\,\mathrm dr\)。因此
\[
 u\in L^2(B),\quad |\nabla u|\in L^2(B),\quad
 u^2\in L^2(B),\quad |\nabla(u^2)|\in L^1(B).
\]
这里所用的奇性指数依次为 \(2/3\)、\(8/3\)、\(4/3\)、\(5/3\)，都小于 \(3\)。

还须确认经典梯度的候选确实是弱梯度，不能只凭原点外的计算略去这一点。对任一坐标方向，除了横向参数为零、因而穿过原点的那一条直线外，函数 \(u\) 和 \(u^2\) 在相应球内截线上都光滑，并在每个紧子区间上绝对连续。异常直线的参数集在 \(\mathbb R^2\) 中为零测集。因此两函数都具有 ACL 性质。结合已算出的函数和候选偏导数的可积性，\cref{b1:thm:sobolev-acl}的反向判据给出 \(u\in H^1(B)\) 与 \(u^2\in W^{1,1}(B)\)。

然而
\[
 \int_B|\nabla(u^2)|^2\,\mathrm dx
 =\frac{16\pi}{9}\int_0^1r^{-4/3}\,\mathrm dr=\infty.
\]
若 \(u^2\) 另有一个 \(L^2\) 弱梯度，由弱导数唯一性，它必须与已经得到的 \(L^1\) 弱梯度几乎处处相同，这与上式矛盾。因此 \(u^2\notin H^1(B)\)。

函数 \(u\) 并不本性有界。乘积的弱导数依旧满足 \(\nabla(u^2)=2u\nabla u\)，但两个 \(L^2\) 因子的乘积由 Hölder 不等式只能统一保证属于 \(L^1\)，不能保证仍属于 \(L^2\)。失去的是原有可积指数，而不是乘积求导公式本身。
\end{solution}

\begin{exercise}\label{b1:ex:21-positive-part-continuity}
设 \(\Omega\subset\mathbb R^d\) 为开集、\(1\le p<\infty\)，实值函数
\(u_n\to u\)、\(v_n\to v\) 于 \(W^{1,p}(\Omega)\)。利用正文的正部强连续性证明
\[
 \max\{u_n,v_n\}\longrightarrow\max\{u,v\},\qquad
 \min\{u_n,v_n\}\longrightarrow\min\{u,v\}
\]
于同一 Sobolev 空间。再在 \((-1,1)\) 上用 \(u_n(x)=x-1/n\)、\(u(x)=x\)，说明正部映射在 \(W^{1,\infty}\) 中不连续。
\end{exercise}
% 来源：自拟；由正文正部强连续性推出上下格运算的强连续性，并保留无穷端点反例。
\begin{solution}
差函数满足 \(u_n-v_n\to u-v\) 于 \(W^{1,p}(\Omega)\)。由\cref{b1:prop:sobolev-lattice-strong-continuity}及恒等式
\[
 \max\{a,b\}=b+(a-b)_+,\qquad
 \min\{a,b\}=a-(a-b)_+,
\]
分别把右端的加法、减法和正部运算通过极限，便得到两个结论。

无穷指数时，给定的 \(u_n-u=-1/n\)，且二者一阶导数都为 \(1\)，故 \(u_n\to u\) 于 \(W^{1,\infty}\bigl((-1,1)\bigr)\)。但在 \((0,1/n)\) 上，\(((u_n)_+)'=0\)、\((u_+)'=1\)，所以
\[
 \|((u_n)_+)'-(u_+)'\|_\infty=1.
\]
区间虽不断缩小，本性上确界仍能检测其中的完整跳差。这正是有限 \(p\) 的控制收敛论证不能照搬的地方。
\end{solution}

\begin{optionalexercise}\label{b1:ex:21-coordinate-change-boundaries}
\begin{enumerate}
\item\label{b1:item:21-inverse-lipschitz-required}
映射 \(\Phi(x)=x^3\) 是 \(I=(-1,1)\) 到自身的光滑 Lipschitz 双射。取坐标函数 \(u(x)=x\)，证明 \(u\in H^1(I)\)，但 \(u\circ\Phi^{-1}\notin H^1(I)\)。
\item\label{b1:item:21-bilipschitz-not-second-order}
令 \(\Psi(x)=x+\tfrac12|x|\)。证明它是 \(I=(-1,1)\) 到 \(J=(-1/2,3/2)\) 的双 Lipschitz 映射，却存在一个函数 \(v\)，它属于所有 \(W^{2,p}(J)\)、\(1\le p\le\infty\)，但 \(v\circ\Psi\) 不属于任何 \(W^{2,p}(I)\)。
\end{enumerate}
\end{optionalexercise}
% 来源：自拟；用退化坐标与折线坐标分别说明逆Lipschitz条件和高阶坐标正则性不能省略。
\begin{solution}
第一问中，\(|\Phi'(x)|=3x^2\le3\)，所以 \(\Phi\) 为 Lipschitz 映射；严格单调性及端点极限说明它是题中双射。其逆映射在零点附近并非 Lipschitz，因为
\[
 \frac{|\Phi^{-1}(h)-\Phi^{-1}(0)|}{|h|}=|h|^{-2/3}\longrightarrow\infty.
\]
坐标函数 \(u\) 及其一阶导数都平方可积。复合函数则为
\[
 (u\circ\Phi^{-1})(x)=\operatorname{sgn}(x)|x|^{1/3}.
\]
它在原点外的经典导数是 \(\tfrac13|x|^{-2/3}\)。若有平方可积的弱导数，该弱导数在两个不含原点的开区间上必须与经典导数几乎处处相同，但
\[
 \int_0^1\left|\frac13x^{-2/3}\right|^2\,\mathrm dx=\infty,
\]
所以不存在这样的弱导数。正向映射光滑且 Lipschitz，并不足以控制经过其逆映射复合后的 Sobolev 范数。

第二问中，\(\Psi\) 在负半轴的斜率为 \(1/2\)，在正半轴的斜率为 \(3/2\)。对任意 \(x<y\)，分段积分给出
\[
 \frac12(y-x)\le\Psi(y)-\Psi(x)\le\frac32(y-x).
\]
故它及其逆映射都 Lipschitz，端点极限给出像恰为 \(J\)。

取 \(v(y)=y\)。在有界区间 \(J\) 上，\(v\)、\(v'=1\)、\(v''=0\) 属于所有所列 \(L^p\) 空间，因此 \(v\in W^{2,p}(J)\)。复合 \(v\circ\Psi=\Psi\) 的弱一阶导数为
\[
 \Psi'=\frac12\mathbf1_{(-1,0)}+\frac32\mathbf1_{(0,1)}.
\]
这个导数在原点的跳跃为 \(1\)。由\cref{b1:ex:20-piecewise-jump-derivatives}的跳跃公式，或直接在原点两侧分部积分，
\[
 \partial^2T_\Psi=\delta_0.
\]
点质量不是正则分布，因此 \(\Psi\) 没有局部可积的二阶弱导数，更不属于任何所列 \(W^{2,p}(I)\)。双 Lipschitz 条件适用于本章的一阶坐标变换，不能自动提升为任意阶的结论。
\end{solution}

\begin{exercise}\label{b1:ex:21-acl-representative-choices}
在 \(\Omega=(-1,1)^2\) 上考虑零函数以及
\[
 v(x,y)=\mathbf1_{\mathbb Q}(x),\quad
 w(x,y)=\mathbf1_{\{(0,0)\}}(x,y).
\]
证明三者属于同一个 \(W^{1,p}(\Omega)\) 等价类，\(1\le p\le\infty\)，但 \(v\) 在每条水平线上都不是绝对连续函数，\(w\) 却仍具有同时针对两个坐标方向的 ACL 性质。解释为何这既不违背共同 ACL 代表的存在性，也不违背一维绝对连续代表的唯一性。最后说明两次零集改值不改变几乎处处的近似微分。
\end{exercise}
% 来源：自拟；从零集改值推进到ACL的逐方向量词及高维代表非逐点唯一性，不重复正文普通可微反例。
\begin{solution}
集合 \(\mathbb Q\cap(-1,1)\) 为一维零测集，由 Fubini 定理，\(v\) 的非零集合在 \(\Omega\) 中为二维零测集；\(w\) 的非零集合也为零测集。因此两者都几乎处处等于零，定义中的函数积分与弱导数积分均不变。它们代表的是零 Sobolev 元素，所有弱偏导数也都是零等价类。

固定任意 \(y\in(-1,1)\)，函数 \(x\mapsto v(x,y)\) 是有理点的指示函数。每个区间都同时含有有理点与无理点，故该截面处处不连续，特别不绝对连续。因此 \(v\) 连“几乎每条水平线”上的绝对连续性都不满足。它在竖直线上为常数，并不能补救另一个坐标方向的失败。

对 \(w\)，若 \(y\ne0\)，则整个水平截面恒为零；若 \(x\ne0\)，则整个竖直截面也恒为零。每个方向只须排除参数为零的那一条直线。因此 \(w\) 是一个同时满足两个方向 ACL 条件的代表，却在原点不等于处处为零的代表。共同 ACL 代表的存在性并没有说每个任意改值的代表都可以使用，也没有断言高维 ACL 代表处处唯一。

一维绝对连续代表的唯一性没有失效。在原点所在的那两条坐标线上，\(w\) 的截面并不连续，因而它们正是 ACL 条件允许排除的异常截面，不能拿来与零截面应用一维唯一性。在其余合格截线上，两者确实处处相同。

最后设 \(E\) 为某一次改值的二维零测集，取 \(z\notin E\)。改值后的函数在 \(z\) 取值为零，在 \(E\) 外也处处为零。对任意 \(\eta>0\)，
\[
 \{\,h:0<|h|<r,\ |f(z+h)-f(z)|>\eta|h|\,\}
 \subseteq E-z.
\]
右侧在每个球中都为零测集，所以左侧的相对测度为零。这就说明近似微分在 \(z\) 存在且为零。仅有改值集上的基点可能不满足这个论证，而它本身是零测集，故几乎处处近似微分保持不变。
\end{solution}

\begin{optionalexercise}\label{b1:ex:21-highest-derivative-kernel}
设 \(\Omega\subset\mathbb R^d\) 为非空连通开集，整数 \(k\ge1\)，\(1\le p\le\infty\)，且 \(u\in W^{k,p}_{\mathrm{loc}}(\Omega)\)。证明以下两条等价：
\begin{equivalences}
\item\label{b1:item:21-highest-derivatives-zero}
所有 \(|\alpha|=k\) 阶弱导数 \(\partial^\alpha u\) 都为零。
\item\label{b1:item:21-polynomial-representative}
函数 \(u\) 几乎处处等于一个总次数不超过 \(k-1\) 的多项式。
\end{equivalences}
去掉连通性后，应怎样修改第二条？
\end{optionalexercise}
% 来源：自拟；局部磨光、有限维闭性与连通拼合的综合应用，不预借Poincare不等式或全域稠密定理。
\begin{solution}
多项式的经典导数也是弱导数，所以第二条直接蕴涵第一条。现假设第一条成立。

固定开球 \(B=B(x_0,r)\)，使 \(\overline{B(x_0,2r)}\subset\Omega\)。由分布卷积的微分公式，当 \(\varepsilon>0\) 足够小时，局部磨光 \(u_\varepsilon=u*\rho_\varepsilon\) 在 \(B\) 的邻域光滑，并满足
\[
 \partial^\alpha u_\varepsilon=0\quad(|\alpha|=k).
\]
对 \(x\in B\)，沿连接 \(x_0\) 与 \(x\) 的线段令 \(F(t)=u_\varepsilon\bigl(x_0+t(x-x_0)\bigr)\)。链式求导表明 \(F^{(k)}=0\)，因为展开后的每项都含一个 \(k\) 阶偏导数。由一维 Taylor 公式，
\[
 u_\varepsilon(x)
 =\sum_{|\alpha|<k}
     \frac{\partial^\alpha u_\varepsilon(x_0)}{\alpha!}(x-x_0)^\alpha.
\]
因此每个 \(u_\varepsilon|_B\) 都是总次数不超过 \(k-1\) 的多项式的限制。

函数 \(u\) 在 \(\overline B\) 的邻域局部可积，即使 \(p=\infty\) 也如此。由\cref{b1:prop:local-smoothing-approximation}的 \(L^1\) 版本，\(u_\varepsilon\to u\) 于 \(L^1(B)\)。固定次数多项式在 \(B\) 上的限制构成一个有限维线性子空间；其 \(L^1(B)\) 范数确实是范数，因为连续多项式若几乎处处为零，便在开球上处处为零，从而是零多项式。由有限维子空间的闭性，\(u|_B\) 几乎处处等于某个次数不超过 \(k-1\) 的多项式 \(P_B\)。这里所用的逼近只有局部 \(L^1\) 收敛，不需要无穷范数的磨光收敛。

若两个这样的球 \(B_1,B_2\) 相交，则 \(P_{B_1}-P_{B_2}\) 在非空开集 \(B_1\cap B_2\) 上几乎处处为零。连续性使其在那里处处为零，多项式恒等性进而给出 \(P_{B_1}=P_{B_2}\) 作为全空间多项式。为看出后一事实，可在交集内取一个小球；多项式在那里为零，其在球心的全部偏导数也为零，而有限 Taylor 展开恰是这个多项式本身。

上述小球覆盖 \(\Omega\)。固定其中一个球，取所有能够通过有限条相交小球链与它连接的球的并集。这个并集开，补集也是其余链类的球的并集，因而也开。由 \(\Omega\) 连通，只能有一个链类，所以全部 \(P_B\) 都是同一个多项式 \(P\)。再由开覆盖的可数子覆盖性，球内各自的零测例外可以并成一个零测集，从而 \(u=P\) 几乎处处于整个 \(\Omega\)。

若 \(\Omega\) 不连通，则在每个连通分支上分别得到一个总次数不超过 \(k-1\) 的多项式，不同分支间不要求系数相同。反过来，按分支给出的这类函数具有所需的局部正则性。确实，若 \(K\Subset\Omega\)，则 \(\delta=\operatorname{dist}(K,\Omega^c)>0\)。对每个 \(x\in K\)，连通球 \(B(x,\delta/2)\) 包含于 \(\Omega\)，因而完全落在 \(x\) 所属的连通分支中；由紧性可取有限个这样的球覆盖 \(K\)，所以 \(K\) 只遇到有限个分支。分支边界不在 \(\Omega\) 内，也不会产生内部跳跃项。
\end{solution}
